\documentclass[11pt]{article}
\usepackage[margin=1in]{geometry}
\usepackage{amsmath, amssymb, amsthm}
\usepackage{mathtools}
\usepackage{xcolor}
\usepackage{hyperref}
\usepackage{mathrsfs}
\usepackage{booktabs}
\usepackage{algorithm}
\usepackage{algpseudocode}
\title{Merged Proofs of RMA}
\author{}
\date{}
\begin{document}
\maketitle
\tableofcontents


% ===== Begin q1_solution.tex =====
\clearpage
\section*{Problem 1}
\addcontentsline{toc}{section}{Problem 1}
\begingroup
\hypersetup{colorlinks=true,
            linkcolor=blue!50!black,
            citecolor=blue!50!black,
            urlcolor=blue!50!black}
\newcommand{\qoneT}{\mathbb{T}}
\newcommand{\qoneR}{\mathbb{R}}
\newcommand{\qoneN}{\mathbb{N}}
\newcommand{\qoneC}{\mathcal{C}}
\newcommand{\qoneD}{\mathcal{D}}
\newcommand{\qoneE}{\mathbb{E}}
\newcommand{\qoneHp}{H}      % Hermite
\newcommand{\qonePN}{P_N}
\newcommand{\qoneLin}{\mathbf{l}}   % stationary linear solution (placeholder for tree symbol)
\newcommand{\qoneLsq}{\mathbf{q}}   % Wick square
\newcommand{\qoneLcb}{\mathbf{c}}   % Wick cube
\newcommand{\qoneIqc}{\mathbf{C}}   % solution driven by the cube
\newtheorem{qonetheorem}{Theorem}
\newtheorem{qonelemma}{Lemma}
\newtheorem{qoneproposition}{Proposition}
\newtheorem{qonecorollary}{Corollary}
\title{Solution to Problem 1: (Lack of) Quasi-Shift-Invariance\\
of the $\Phi^4_3$ Measure}
\author{}
\date{}
\maketitle


\section*{Statement and Answer}

Let $\qoneT^3$ be the unit $3$-torus and $\mu$ the $\Phi^4_3$ measure on
$\qoneD'(\qoneT^3)$. For smooth $\psi\not\equiv0$ let $T_\psi(u)=u+\psi$, and let
$T_\psi^*\mu$ be the pushforward.

\medskip
\noindent\textbf{Theorem.} \emph{For every smooth $\psi\not\equiv0$ and every
admissible choice of mass and (nonzero) coupling constant, the measures $\mu$ and
$T_\psi^*\mu$ are mutually \textbf{singular}.}

\medskip
\noindent The answer is \textbf{NO}: $\mu$ and $T_\psi^*\mu$ are \emph{not}
equivalent. This is a simplified form of a result of Hairer--Peroni and rests on
ideas of Hairer--Kusuoka--Nagoji.

\medskip
\noindent\textbf{Warning on a fatal false premise.} It is tempting to write
$d\mu/d\mu_0=Z^{-1}e^{-\mathcal R(u)}$ against the Gaussian free field $\mu_0$
and deduce quasi-invariance by a Cameron--Martin/Wick-binomial computation. This
is \emph{wrong in three dimensions}: Glimm showed that, unlike $d=2$, the
$\Phi^4_3$ measure $\mu$ and the free field are mutually \emph{singular} in
$d=3$ (the borderline dimension is $8/3$). There is no such Radon--Nikodym
density, so any argument built on it collapses. The correct proof exhibits an
event separating $\mu$ from $T_\psi^*\mu$, driven by a genuinely
$3$-dimensional logarithmic divergence.

\section*{Notation}

Fix a space--time white noise $\xi$ on $\qoneR\times\qoneT^3$ and let $\qoneLin$ be the
stationary solution of $(\partial_t+1-\Delta)\qoneLin=\xi$. Write $\qonePN=P_{\le N}$ for
the projection onto Fourier modes $|k|\le N$, $\qoneLin_N=\qonePN\qoneLin$, and
$c_N=\qoneE\qoneLin_N^2$. Let $\qoneHp_n$ be the Hermite polynomials normalized by
$\qoneHp_0\equiv1$, $\qoneHp_n'=n\qoneHp_{n-1}$, $\qoneE\qoneHp_n(Z,1)=0$ for standard normal $Z$. Set
\[
\qoneLsq=\lim_{N}\qoneHp_2(\qoneLin_N,c_N),\qquad \qoneLcb=\lim_N\qoneHp_3(\qoneLin_N,c_N),
\]
the Wick square and cube (convergent in $\qoneC^{-1-2\kappa}$ and $\qoneC^{-3/2-3\kappa}$
respectively), and let $\qoneIqc$ solve $(\partial_t+1-\Delta)\qoneIqc=\qoneLcb$. We fix
$0<\kappa\le 1/100$. By \cite{HKN24}, $\qoneLcb$ and $\qoneIqc$ are a.s.\ continuous in
time with values in $\qoneC^{-1/2-\kappa}$ and $\qoneC^{1/2-3\kappa}$.

We will need the additional constant
\begin{equation}\label{eq:cN2}
c_{N,2}:=\qoneE\bigl[\qoneLsq_N\,\qoneIqc_N\bigr],\qquad \qoneLsq_N:=\qonePN\qoneLsq,\ \qoneIqc_N:=\qonePN\qoneIqc.
\end{equation}

\section*{The decomposition of $\mu$}

\begin{qoneproposition}[\cite{HM18,EW24,CC18}]\label{prop:decomp}
There is a stationary process $v$, a.s.\ continuous in time with values in
$\qoneC^{1-2\kappa}(\qoneT^3)$, such that
\begin{equation}\label{eq:u}
u=\qoneLin-\qoneIqc+v
\end{equation}
is stationary with fixed-time law $\mu$. Moreover $v$ admits the paracontrolled
form
\begin{equation}\label{eq:v}
v=-3\bigl((v-\qoneIqc)\prec\qoneLin\bigr)+v^\sharp,
\end{equation}
where $v^\sharp\in\qoneC^{1+4\kappa}$ and $\prec$ is Bony's paraproduct.
\end{qoneproposition}

\section*{The distinguishing event}

For $\gamma>0$ define
\[
B_\gamma:=\Bigl\{u\in\qoneD':\ \lim_{N\to\infty}\bigl\langle(\log N)^{-\gamma}\bigl(\qoneHp_3(\qonePN u;c_N)+9c_{N,2}\qonePN u\bigr),\psi\bigr\rangle_{\qoneT^3}=0\Bigr\},
\]
where $N$ runs over $2^{\qoneN}$. The next two lemmas (proved as in
\cite[Lems.~3.11--3.12]{HKN24}) supply the analytic input.

\begin{qonelemma}\label{lem:cube}
For $\gamma>\tfrac12$ and each fixed $t>0$,
$(\log N)^{-\gamma}\,{:}\qoneLin_N^3{:}(t)\to0$ a.s.\ in $\qoneC^{-3/2}(\qoneT^3)$ and in
$L^p(\Omega;\qoneC^{-3/2})$ for all $p>0$.
\end{qonelemma}

\begin{proof}[Sketch]
With $\beta=-d/2=-3/2$ and the embedding $W^{\beta,2p}\hookrightarrow
\qoneC^{\beta-d/2p}$, one bounds
$\qoneE\|(\log N)^{-\gamma}{:}\qoneLin_N^3{:}\|_{W^{-3/2,2p}}^{2p}$ by a constant times
\[
(\log N)^{-2\gamma}\sum_{|\omega_i|\le N}\langle\omega_1+\omega_2+\omega_3\rangle^{-3}\prod_{i}\langle\omega_i\rangle^{-2}
\ \lesssim\ (\log N)^{-2\gamma+1},
\]
which is summable along $N\in2^\qoneN$ once $\gamma>\tfrac12$; Borel--Cantelli
concludes.
\end{proof}

\begin{qonelemma}\label{lem:logdiv}
For $N$ large, $c_{N,2}\gtrsim\log N$.
\end{qonelemma}

\begin{proof}[Sketch]
Expanding \eqref{eq:cN2} and integrating the heat kernels in time,
\[
c_{N,2}\gtrsim\sum_{|\omega_i|\le N}\frac{1}{\langle\omega_1\rangle^2\langle\omega_2\rangle^2}\cdot\frac{1}{\langle\omega_1\rangle^2+\langle\omega_2\rangle^2+\langle\omega_1+\omega_2\rangle^2}
\gtrsim\sum_{|\omega_1|\le|\omega_2|\le N}\frac{1}{(1+|\omega_1|^2)(1+|\omega_2|^4)}.
\]
Bounding the sum below by an integral gives
$\int_0^N\frac{r}{1+r^2}\,dr\simeq\log N$.
\end{proof}

The following standard facts are recorded for reference.

\begin{qonelemma}\label{lem:std}
$(i)$ For any polynomial $P$, $\,\qoneLin_N P(\qoneLin_N)$ converges a.s.\ in
$\qoneC^{-1/2-\kappa}$. $(ii)$ The expressions ${:}\qoneLin_N^2{:}\,\qoneIqc_N-3c_{N,2}\qoneIqc_N$
and ${:}\qoneLin_N^2{:}\,v_N+3c_{N,2}(v_N-\qoneIqc_N)$ both converge a.s.\ in
$\qoneC^{-1-2\kappa}$ as $N\to\infty$.
\end{qonelemma}

\section*{Step 1: $\mu(B_\gamma)=1$}

Fix $\gamma>\tfrac12$. Let $u$ be as in \eqref{eq:u} and write $u_N=\qonePN u$.
Expanding the Wick cube and using $u_N=\qoneLin_N-\qoneIqc_N+v_N$,
\begin{align*}
(\log N)^{-\gamma}\qoneHp_3(u_N;c_N)
&=(\log N)^{-\gamma}{:}\qoneLin_N^3{:}
-3(\log N)^{-\gamma}{:}\qoneLin_N^2{:}\,\qoneIqc_N\\
&\quad+3(\log N)^{-\gamma}{:}\qoneLin_N^2{:}\,v_N
+(\log N)^{-\gamma}R_N,
\end{align*}
where $R_N$ collects the remaining mixed terms. By Lemma~\ref{lem:cube} the
first term $\to0$; by Lemma~\ref{lem:std}(i) and standard product estimates the
terms in $R_N$ $\to0$. It remains to treat
$-{:}\qoneLin_N^2{:}\,\qoneIqc_N+{:}\qoneLin_N^2{:}\,v_N+3c_{N,2}u_N$, which we regroup as
\[
\Bigl({:}\qoneLin_N^2{:}\,v_N+3c_{N,2}(v_N-\qoneIqc_N)\Bigr)-\Bigl({:}\qoneLin_N^2{:}\,\qoneIqc_N-3c_{N,2}\qoneIqc_N\Bigr).
\]
By Lemma~\ref{lem:std}(ii) this converges a.s.\ in $\qoneC^{-1-2\kappa}$; multiplied
by $(\log N)^{-\gamma}\to0$ it tends to $0$. Pairing with $\psi$,
\[
\bigl\langle(\log N)^{-\gamma}\bigl(\qoneHp_3(u_N;c_N)+9c_{N,2}u_N\bigr),\psi\bigr\rangle\xrightarrow[N\to\infty]{}0\quad\text{a.s.},
\]
i.e.\ $u(t)\in B_\gamma$ for fixed $t$. Hence $\mu(B_\gamma)=1$.

\section*{Step 2: $(T_\psi^*\mu)(B_\gamma)=0$}

We must show $u+\psi\notin B_\gamma$. Expanding with $\psi_N=\qonePN\psi$,
\begin{align*}
&(\log N)^{-\gamma}\qoneHp_3(u_N+\psi_N;c_N)+(\log N)^{-\gamma}9c_{N,2}(u_N+\psi_N)\\
&\quad=\Bigl[(\log N)^{-\gamma}\qoneHp_3(u_N;c_N)+9(\log N)^{-\gamma}c_{N,2}u_N\Bigr]\\
&\qquad+3(\log N)^{-\gamma}{:}u_N^2{:}\,\psi_N+3(\log N)^{-\gamma}u_N\psi_N^2\\
&\qquad+(\log N)^{-\gamma}\psi_N^3+9(\log N)^{-\gamma}c_{N,2}\psi_N.
\end{align*}
The bracketed term tends to $0$ a.s.\ (Step 1). Since ${:}u_N^2{:}$ and $u_N$
converge to finite distributional limits (Lemma~\ref{lem:std}(i)), the next three
terms tend to $0$ after multiplication by $(\log N)^{-\gamma}$. The decisive term
is the last: by Lemma~\ref{lem:logdiv},
\[
9(\log N)^{-\gamma}c_{N,2}\,\langle\psi_N,\psi\rangle\ \gtrsim\ (\log N)^{-\gamma+1}\,\langle\psi_N,\psi\rangle .
\]
As $N\to\infty$, $\langle\psi_N,\psi\rangle\to\|\psi\|_{L^2}^2>0$ (because
$\psi\not\equiv0$), and $(\log N)^{-\gamma+1}\to\infty$ for $\gamma<1$. Choosing
$\tfrac12<\gamma<1$, the displayed quantity \emph{diverges}, so the full pairing
does not converge to $0$; that is, $u+\psi\notin B_\gamma$ a.s. Therefore
$(T_\psi^*\mu)(B_\gamma)=\mu(\{u:u+\psi\in B_\gamma\})=0$.

\section*{Conclusion}

By Steps 1 and 2, $\mu(B_\gamma)=1$ while $(T_\psi^*\mu)(B_\gamma)=0$. Thus $\mu$
and $T_\psi^*\mu$ are concentrated on disjoint Borel sets and are mutually
singular. The only structural inputs were that $\psi\not\equiv0$ (used solely to
ensure $\|\psi\|_{L^2}^2>0$) and that the coupling is nonzero, so the conclusion
holds for all such $\psi$. \qed

\section*{Remarks}

\begin{enumerate}\setlength\itemsep{4pt}
\item \textbf{Where dimension three enters.} The entire effect is carried by
$c_{N,2}=\qoneE[\qoneLsq_N\qoneIqc_N]\gtrsim\log N$ (Lemma~\ref{lem:logdiv}). This logarithmic
divergence is absent in $d=2$, where $\mu$ \emph{is} quasi-invariant under smooth
shifts and even equivalent to the free field. The borderline dimension for
shift-quasi-invariance is exactly $3$, and the question is on which side it falls;
the answer is the singular side.

\item \textbf{Why the heuristic ``density'' argument fails.} Formally
$dT_\psi^*\mu/d\mu$ involves a term proportional to the would-be field
$\qoneLin^3-C\qoneLin$, which does \emph{not} define a random distribution in $d=3$ (its
covariance behaves like $|x-y|^{-3}$, non-integrable). The log-divergent constant
in the formal density is precisely what makes the candidate density blow up. The
rigorous proof above replaces the nonexistent density by the cut-off functional
defining $B_\gamma$ and extracts the divergence cleanly.

\item \textbf{Role of $\gamma$.} The window $\tfrac12<\gamma<1$ is forced:
$\gamma>\tfrac12$ kills the cube and the chaos remainders (Lemma~\ref{lem:cube}),
while $\gamma<1$ lets the shift term $(\log N)^{-\gamma}c_{N,2}\asymp(\log
N)^{1-\gamma}$ diverge. Any $\gamma$ in this range distinguishes the measures.
\end{enumerate}

\begin{thebibliography}{9}
\bibitem{HP} M.~Hairer, J.~Peroni.
  \emph{Quasi shift invariance of $\Phi^4$ measures}. To appear.
\bibitem{HKN24} M.~Hairer, S.~Kusuoka, H.~Nagoji.
  \emph{Singularity of solutions to singular SPDEs}.
  Preprint, arXiv:2409.10037, 2024.
\bibitem{Glimm68} J.~Glimm.
  \emph{Boson fields with the $\Phi^4$ interaction in three dimensions}.
  Comm.\ Math.\ Phys.\ \textbf{10} (1968), 1--47.
\bibitem{HM18} M.~Hairer, K.~Matetski.
  \emph{Discretisations of rough stochastic PDEs}.
  Ann.\ Probab.\ \textbf{46} (2018), 1651--1709.
\bibitem{EW24} S.~Esquivel, H.~Weber.
  \emph{A priori bounds for the dynamic fractional $\Phi^4$ model on $\qoneT^3$}.
  Preprint, arXiv:2411.16536, 2024.
\bibitem{CC18} R.~Catellier, K.~Chouk.
  \emph{Paracontrolled distributions and the $3$-dimensional stochastic
  quantization equation}. Ann.\ Probab.\ \textbf{46} (2018), 2621--2679.
\bibitem{Hairer14} M.~Hairer.
  \emph{A theory of regularity structures}. Invent.\ Math.\ \textbf{198} (2014),
  269--504.
\end{thebibliography}

\endgroup
% ===== End q1_solution.tex =====

% ===== Begin q2_solution.tex =====
\clearpage
\section*{Problem 2}
\addcontentsline{toc}{section}{Problem 2}
\begingroup
\hypersetup{colorlinks=true,
            linkcolor=blue!50!black,
            citecolor=blue!50!black,
            urlcolor=blue!50!black}
\newcommand{\qtwoF}{F}
\newcommand{\qtwooo}{\mathfrak{o}}
\newcommand{\qtwopp}{\mathfrak{p}}
\newcommand{\qtwoqq}{\mathfrak{q}}
\newcommand{\qtwoC}{\mathbb{C}}
\newcommand{\qtwoZ}{\mathbb{Z}}
\newcommand{\qtwoGL}{\mathrm{GL}}
\newcommand{\qtwoMn}{M_n}
\newcommand{\qtwoSe}{\mathcal{S}_{\mathrm e}}
\newcommand{\qtwoSw}{\mathcal{S}}
\newcommand{\qtwoW}{\mathcal{W}}
\newcommand{\qtwovol}{\operatorname{vol}}
\newcommand{\qtwodiag}{\operatorname{diag}}
\newcommand{\qtwotr}{\operatorname{trace}}
\newcommand{\qtwottilde}[1]{\widetilde{#1}}
\newcommand{\qtwolRS}{\ell_{\mathrm{RS}}}
\newtheorem{qtwotheorem}{Theorem}
\newtheorem{qtwolemma}{Lemma}
\newtheorem{qtwoproposition}{Proposition}
\newtheorem{qtwocorollary}{Corollary}
\title{Solution to Problem 2: A Universal Test Vector\\
for Rankin--Selberg Integrals}
\author{}
\date{}
\maketitle


\section*{Statement and Answer}

Let $\qtwoF$ be a non-archimedean local field with ring of integers $\qtwooo$, maximal
ideal $\qtwopp$, residue cardinality $q$, and a nontrivial additive character
$\psi:\qtwoF\to\qtwoC^\times$ of conductor $\qtwooo$. Write $G_r:=\qtwoGL_r(\qtwoF)$, let $N_r<G_r$
be the upper-triangular unipotent subgroup, and embed $G_n\hookrightarrow
G_{n+1}$ as the upper-left block. Set $K_r:=\qtwoGL_r(\qtwooo)$, with Haar measures
giving $K_r$ and $N_r\cap K_r$ volume one.

\medskip
\noindent\textbf{Theorem.} \emph{Let $\Pi$ be a generic irreducible admissible
representation of $G_{n+1}$, realized in its $\psi^{-1}$-Whittaker model
$\qtwoW(\Pi,\psi^{-1})$. Then there exists a single $W\in\qtwoW(\Pi,\psi^{-1})$ with the
following property. For every generic irreducible admissible $\pi$ of $G_n$ with
conductor ideal $\qtwoqq$, generator $Q$ of $\qtwoqq^{-1}$, and
$u_Q:=I_{n+1}+QE_{n,n+1}$, there exists $V\in\qtwoW(\pi,\psi)$ such that}
\[
\int_{N_n\backslash G_n} W(\qtwodiag(g,1)u_Q)\,V(g)\,|\det g|^{s-\frac12}\,dg
\]
\emph{is finite and nonzero for all $s\in\qtwoC$.}

\medskip
\noindent The answer is \textbf{YES}. The single $W$ may be taken to be an
explicit, $\Pi$-independent-of-$\pi$ Kirillov-model vector, and for each $\pi$
the role of $V$ is played by the \emph{normalized newvector}. The integral is in
fact \emph{constant in $s$} and equal to an explicit nonzero scalar; this is why
finiteness and nonvanishing hold for all $s$, with no analytic continuation or
pole-cancellation required.

\medskip
\noindent\textbf{A failure mode to avoid.} One cannot choose $V$ as a
Schwartz bump in the Kirillov model making the integrand constant on its
support: $V$ has a central character, so it cannot be compactly supported modulo
the center. (Already for $n=1$ the relevant integral is a normalized Gauss sum,
whose integrand is genuinely non-constant.) The correct argument runs through the
Godement--Jacquet functional equation, not through test-vector flexibility.

\section*{Choice of $W$ and $V$}

For $x\in\qtwoF$ and $y\in\qtwoF^\times$ write
\[
d_y:=\qtwodiag(1,\dots,1,y)\in G_n\hookrightarrow G_{n+1},\qquad
u_x:=I_{n+1}+xE_{n,n+1}\in N_{n+1}.
\]

\noindent\textbf{The vector $W$.} Define $W_0$ on $G_n$ by
\begin{equation}\label{eq:W0}
W_0(g):=\int_{N_n}\mathbf 1_{K_n}(xg)\,\psi(x)\,dx,
\end{equation}
a $\psi^{-1}$-Whittaker function on $G_n$ which extends, via the theory of the
Kirillov model \cite{Bernstein}, to an element of $\qtwoW(\Pi,\psi^{-1})$ on
$G_{n+1}$. We take $W:=u_Q W_0$ (equivalently we evaluate $W_0$ at the shifted
argument); $W_0$ depends only on $\Pi$, not on $\pi$.

\noindent\textbf{The vector $V$.} Given $\pi$ with conductor $\qtwoqq$, let
$V\in\qtwoW(\pi,\psi)$ be the normalized newvector: the unique $K_1(\qtwoqq)$-invariant
vector with $V(1)=1$, where
$K_1(\qtwoqq)=\{g\in K_n:\text{last row}\equiv(0,\dots,0,1)\bmod\qtwoqq\}$
\cite{JPSS81,Matringe}.

We choose a generator $Q$ of $\qtwoqq^{-1}$, so $|Q|=[\qtwooo:\qtwoqq]$. Recall the
conductor is characterized by the $\varepsilon$-factor:
\begin{equation}\label{eq:eps}
\varepsilon(\tfrac12+s,\pi,\psi)=|Q|^{-s}\,\varepsilon(\tfrac12,\pi,\psi).
\end{equation}

\section*{The Godement--Jacquet engine}

\begin{qtwolemma}[Godement--Jacquet functional equation \cite{GJ72}]\label{lem:GJ}
Let $f$ be a matrix coefficient of $\pi$ and $\phi\in\qtwoSw(\qtwoMn(\qtwoF))$. The zeta
integral
\[
Z(\phi,f,s):=\int_{G_n}\phi(g)f(g)\,|\det g|^{\frac{n-1}{2}+s}\,dg
\]
converges for $\Re(s)$ large, continues meromorphically with
$Z(\phi,f,s)/L(s,\pi)$ holomorphic, and satisfies
\[
\gamma(s,\pi,\psi)\,Z(\phi,f,s)=Z(\widehat\phi,f^\vee,1-s),\qquad
\gamma(s,\pi,\psi)=\varepsilon(s,\pi,\psi)\frac{L(1-s,\qtwottilde\pi)}{L(s,\pi)},
\]
where $f^\vee(g)=f(g^{-1})$ and $\widehat\phi$ is the $\psi$-self-dual Fourier
transform on $\qtwoMn(\qtwoF)$. The identities persist when $f$ is taken to be a
Whittaker function of $\pi$.
\end{qtwolemma}

Let $\qtwoSe(\qtwoF^\times)$ be the Schwartz--Bruhat functions $\beta$ with
$\beta(xy)=\beta(x)$ whenever $|y|=1$ (i.e.\ $\beta$ depends only on $|x|$),
with Mellin inversion $\beta(y)=\int_{(\sigma)}\tilde\beta(s)|y|^s\,ds$. For such
$\beta$ define $\beta^\sharp$ by
$\beta^\sharp(y)=\int_{(\sigma)}\tilde\beta(s)|y|^{-s}\,ds/\gamma(\tfrac12+s,\pi,\psi)$.

\begin{qtwolemma}\label{lem:beta}
Define $\beta$ by $\tilde\beta(s)=\varepsilon(\tfrac12+s,\pi,\psi)/L(\tfrac12+s,\pi)$.
Then:
\begin{enumerate}\setlength\itemsep{2pt}
\item $\beta$ is supported on $\{|Q|\le|y|\le|Q|q^n\}$ and equals
$\varepsilon(\tfrac12,\pi,\psi)$ on $\{|y|=|Q|\}$;
\item $\beta^\sharp$ has Mellin transform $\widetilde{\beta^\sharp}(s)=1/L(\tfrac12+s,\qtwottilde\pi)$,
is supported on $\{1\le|y|\le q^n\}$, and equals $1$ on $\qtwooo^\times=\{|y|=1\}$.
\end{enumerate}
\end{qtwolemma}

\begin{proof}
By definition of $\gamma$, $\widetilde{\beta^\sharp}(s)=\tilde\beta(s)/\gamma(\tfrac12+s,\pi,\psi)
=L(\tfrac12-s,\qtwottilde\pi)^{-1}\cdot(\cdots)$ simplifies to $1/L(\tfrac12+s,\qtwottilde\pi)$.
The inverse $L$-values are monic polynomials in $q^{-s}$ of degree $\le n$, so by
Mellin inversion and \eqref{eq:eps}, $\beta$ and $\beta^\sharp$ have the stated
supports and boundary values.
\end{proof}

\begin{qtwolemma}\label{lem:transfer}
For $\pi$ unitary generic, with $f$ a Whittaker function and
$\phi\in\qtwoSw(\qtwoMn(\qtwoF))$,
\[
\int_{G_n}\phi(g)f(g)\beta(\det g)|\det g|^{n/2}\,dg
=\int_{G_n}\widehat\phi(g)f^\vee(g)\beta^\sharp(\det g)|\det g|^{n/2}\,dg.
\]
\end{qtwolemma}

\begin{proof}
Insert the Mellin expansion of $\beta$ at $\sigma=0$; the resulting double
integral converges absolutely, and is recognized as
$\int_{(0)}\tilde\beta(s)Z(\phi,f,\tfrac12+s)\,ds$. Apply the functional equation
(Lemma~\ref{lem:GJ}) and recognize the result as the right-hand side.
\end{proof}

\section*{The local computation}

Set $t_Q:=d_Q^{-1}u_Q=u_1 d_Q^{-1}$.

\begin{qtwolemma}\label{lem:relations}
There exist $\beta\in\qtwoSe(\qtwoF^\times)$ and $\phi\in\qtwoSw(\qtwoMn(\qtwoF))$ with, for all
$g\in G_n$,
\begin{align}
\int_{N_n}\beta(\det xg)\,\phi(xg)\,\psi(x)\,dx
&=\varepsilon(\tfrac12,\pi,\psi)\,W_0(g\,t_Q),\label{eq:rel1}\\
\beta^\sharp(\det g)\,\widehat\phi(g)&=|Q|^n\,\mathbf 1_{K_1(\qtwoqq)}(g).\label{eq:rel2}
\end{align}
\end{qtwolemma}

\begin{proof}
Take $\beta$ as in Lemma~\ref{lem:beta} and
$\phi(x):=\psi(-x_{nn})\phi_0(xd_Q^{-1})$ with $\phi_0=\mathbf 1_{\qtwoMn(\qtwooo)}$. For
\eqref{eq:rel1}, use $W_0(gt_Q)=\psi(-g_{nn})W_0(gd_Q^{-1})$ and the identities
$\phi(xg)=\psi(-g_{nn})\phi_0(xgd_Q^{-1})$,
$\beta(\det g)\phi_0(gd_Q^{-1})=\varepsilon(\tfrac12,\pi,\psi)\mathbf 1_{K_n}(gd_Q^{-1})$
(the latter from Lemma~\ref{lem:beta}(1) and
$\mathbf 1_{K_n}(g)=\mathbf 1_{\qtwooo^\times}(\det g)\phi_0(g)$), then integrate
against $\psi(x)\,dx$ recalling \eqref{eq:W0}. For \eqref{eq:rel2}, the Fourier
computation $\widehat\phi(x)=|Q|^n\widehat{\phi_0}(d_Q(x-E_{nn}))$ with
$\widehat{\phi_0}=\phi_0$, together with Lemma~\ref{lem:beta}(2), gives
$\beta^\sharp(\det g)\widehat\phi(g)=|Q|^n\mathbf 1_{K_1(\qtwoqq)}(g)$, since
$K_1(\qtwoqq)=\{g:d_Q(g-E_{nn})\in\qtwoMn(\qtwooo),\ \det g\in\qtwooo^\times\}$.
\end{proof}

\section*{Conclusion}

For $W\in\qtwoW(\Pi,\psi^{-1})$, $V\in\qtwoW(\pi,\psi)$, write the Rankin--Selberg
integral
\[
\qtwolRS(s,W,V):=\int_{N_n\backslash G_n}W(\qtwodiag(g,1))\,V(g)\,|\det g|^{s-\frac12}\,dg.
\]

\begin{qtwoproposition}\label{prop:final}
Let $W_0$ be as in \eqref{eq:W0} and $V$ the normalized newvector of $\pi$. Then
for all $s\in\qtwoC$,
\[
\qtwolRS(s,u_Q W_0,\,d_Q V)=c\,|Q|^{-n/2},\qquad
c:=\varepsilon(\tfrac12,\pi,\psi)^{-1}|Q|^n\qtwovol(K_1(\qtwoqq))\asymp 1.
\]
\end{qtwoproposition}

\begin{proof}
A change of variables gives the homogeneity
$\qtwolRS(s,u_Q W_0,d_Q V)=|Q|^{-(s-\frac12)}\qtwolRS(s,t_Q W_0,V)$. Since $W_0$ is
supported on $\det^{-1}(\qtwooo^\times)$, the translate $t_Q W_0$ is supported on
$\det^{-1}(Q\qtwooo^\times)$, so $\qtwolRS(s,t_Q W_0,V)$ is a constant multiple of
$|Q|^s$; it therefore suffices to evaluate at $s=\tfrac{n+1}{2}$. Unfolding and
applying Lemmas~\ref{lem:relations} (i.e.\ \eqref{eq:rel1}), \ref{lem:transfer},
and \eqref{eq:rel2} in turn, with $f(g)=V(g)$,
\begin{align*}
\varepsilon(\tfrac12,\pi,\psi)\,\qtwolRS(\tfrac{n+1}{2},t_Q W_0,V)
&=\varepsilon(\tfrac12,\pi,\psi)\!\int_{N_n\backslash G_n}\!W_0(gt_Q)V(g)|\det g|^{n/2}dg\\
&=\int_{G_n}\phi(g)f(g)\beta(\det g)|\det g|^{n/2}dg\\
&=\int_{G_n}\widehat\phi(g)f^\vee(g)\beta^\sharp(\det g)|\det g|^{n/2}dg\\
&=|Q|^n\!\int_{K_1(\qtwoqq)}\!V(g^{-1})|\det g|^{n/2}dg
=|Q|^n\qtwovol(K_1(\qtwoqq)),
\end{align*}
using $K_1(\qtwoqq)$-invariance, $V(1)=1$, and $|\det g|=1$ on $K_1(\qtwoqq)$. Hence
$\qtwolRS(\tfrac{n+1}{2},t_Q W_0,V)=c$, so $\qtwolRS(s,t_Q W_0,V)=c\,|Q|^{s-\frac{n+1}{2}}$
and $\qtwolRS(s,u_Q W_0,d_Q V)=c\,|Q|^{-n/2}$.
\end{proof}

Since $c\,|Q|^{-n/2}$ is a nonzero scalar independent of $s$, the integral is
finite and nonzero for all $s\in\qtwoC$. The vector $W:=u_Q W_0$ (built from the
single, $\pi$-independent $W_0$) and the $\pi$-newvector $V$ (rescaled by $d_Q$)
exhibit the required test vectors. This proves the Theorem. \qed

\section*{Remarks}

The mechanism is transparent: $u_Q W_0$ is supported on a single $\det$-coset, so
the zeta integral collapses to a \emph{constant} in $s$ rather than an
$L$-factor times a unit. The Godement--Jacquet functional equation, packaged
through the $\gamma$-factor transform $\beta\mapsto\beta^\sharp$
(Lemma~\ref{lem:beta}), converts the $\psi$-twisted unipotent shift into an
integral over a translate of $K_1(\qtwoqq)$, where newvector theory yields
$\qtwovol(K_1(\qtwoqq))$. The $n=1$ case recovers the classical normalized Gauss sum and
has appeared in subconvexity and moment applications \cite{Sarnak,MV}; the
general $n$ statement, with a single unipotent translate rather than an average,
is the present lemma (joint work in progress with S.~Jana). A version with an
average over many translates appears in \cite{BKL}.

\begin{thebibliography}{9}
\bibitem{Bernstein} J.~N.~Bernstein.
  \emph{$P$-invariant distributions on $\qtwoGL(N)$ and the classification of
  unitary representations of $\qtwoGL(N)$ (non-archimedean case)}. In
  \emph{Lie group representations II}, LNM \textbf{1041}, Springer, 1984,
  50--102.
\bibitem{BKL} A.~R.~Booker, M.~Krishnamurthy, M.~Lee.
  \emph{Test vectors for Rankin--Selberg $L$-functions}.
  J.\ Number Theory \textbf{209} (2020), 37--48.
\bibitem{GJ72} R.~Godement, H.~Jacquet.
  \emph{Zeta functions of simple algebras}. LNM \textbf{260}, Springer, 1972.
\bibitem{JPSS81} H.~Jacquet, I.~I.~Piatetski-Shapiro, J.~Shalika.
  \emph{Conducteur des repr\'esentations du groupe lin\'eaire}.
  Math.\ Ann.\ \textbf{256} (1981), 199--214.
\bibitem{Matringe} N.~Matringe.
  \emph{Essential Whittaker functions for $\qtwoGL(n)$}.
  Doc.\ Math.\ \textbf{18} (2013), 1191--1214.
\bibitem{MV} P.~Michel, A.~Venkatesh.
  \emph{The subconvexity problem for $\qtwoGL_2$}.
  Publ.\ Math.\ IH\'ES \textbf{111} (2010), 171--271.
\bibitem{Sarnak} P.~Sarnak.
  \emph{Fourth moments of Gr\"ossencharakteren zeta functions}.
  Comm.\ Pure Appl.\ Math.\ \textbf{38} (1985), 167--178.
\end{thebibliography}

\endgroup
% ===== End q2_solution.tex =====

% ===== Begin q3_solution.tex =====
\clearpage
\section*{Problem 3}
\addcontentsline{toc}{section}{Problem 3}
\begingroup
\hypersetup{colorlinks=true,
            linkcolor=blue!50!black,
            citecolor=blue!50!black,
            urlcolor=blue!50!black}
\newcommand{\qthreeR}{\mathbb{R}}
\newcommand{\qthreeN}{\mathbb{N}}
\newcommand{\qthreeZ}{\mathbb{Z}}
\newcommand{\qthreeSnl}{S_n(\lambda)}
\newcommand{\qthreewt}{\operatorname{wt}}
\newtheorem{qthreetheorem}{Theorem}
\newtheorem{qthreelemma}{Lemma}
\newtheorem{qthreeproposition}{Proposition}
\newtheorem{qthreecorollary}{Corollary}
\theoremstyle{definition}
\newtheorem{qthreedefinition}{Definition}
\newtheorem{qthreeremark}{Remark}
\title{Solution to Problem 3: A Probabilistic Interpretation\\
for Interpolation Macdonald Polynomials}
\author{}
\date{}
\maketitle


\section*{Statement and Answer}

Let $\lambda=(\lambda_1>\cdots>\lambda_n\ge0)$ be a partition with distinct
parts that is moreover \emph{restricted}: it has a unique part of size $0$ and no
part of size $1$. Let $\qthreeSnl$ be the orbit of $\lambda$ under permutation of
parts (equivalently, the compositions that rearrange $\lambda$). Let
$F^*_\mu(x_1,\dots,x_n;q,t)$ and $P^*_\lambda(x_1,\dots,x_n;q,t)$ be the
interpolation ASEP polynomial and interpolation Macdonald polynomial.

\medskip
\noindent\textbf{Theorem.} \emph{There is a nontrivial Markov chain on $\qthreeSnl$ ---
the \emph{interpolation $t$-Push TASEP} --- whose stationary distribution is}
\[
\pi^*_\lambda(\mu)\;=\;\frac{F^*_\mu(x_1,\dots,x_n;\,q=1,\,t)}{P^*_\lambda(x_1,\dots,x_n;\,q=1,\,t)},
\qquad \mu\in\qthreeSnl,
\]
\emph{and whose transition probabilities are not described using the
polynomials $F^*_\mu$.}

\medskip
\noindent The answer is \textbf{YES}. We caution at the outset that the chain is
\emph{not} the ordinary multispecies ASEP: that chain has the ordinary (not
interpolation) Macdonald polynomials as its stationary weights (Ayyer--Martin--Williams),
and substituting it here would solve a different, already-known problem. The
correct object is a $q=1$ \emph{push}-TASEP with an interpolation correction
encoded by a ringing ``bell''.

\section*{Probabilities and notation}

Throughout we work at $q=1$ and write $P^*_\lambda=P^*_\lambda(x;1,t)$, etc.
For $1\le k\le n$ put
\[
p_k:=\frac{t^{-n+1}(1-t)}{x_k-t^{-n+2}},\qquad
q_k:=\frac{(1-t)\,x_k}{x_k-t^{-n+2}} .
\]
For $0<t<1$ and $x_i>t^{-n+1}$ these lie in $[0,1]$ and serve as
skipping/displacement probabilities. We use $e^*_k(x;t)$ for the interpolation
elementary symmetric polynomials and the factorization (valid at $q=1$
\cite{Dolega,BDW25a})
\begin{equation}\label{eq:factor}
P^*_\lambda(x;1,t)=\prod_{1\le i\le\lambda_1}e^*_{\lambda'_i}(x;t),
\end{equation}
$\lambda'$ the conjugate partition.

\section*{The interpolation $t$-Push TASEP}

\begin{qthreedefinition}\label{def:chain}
Fix $\lambda$ with $\lambda_n=0$. States of the chain are configurations of
particles on a ring of $n$ sites labelled by the parts of $\lambda$: state $\eta$
places the particle labelled $\eta_j$ at site $j$ (label $0$ is a vacancy). One
step from $\eta$ proceeds in three stages.

\smallskip
\noindent\textbf{(Step 0) The bell.} The bell at site $j$ rings with probability
\[
P_j=\frac{\prod_{k<j}\bigl(x_k-t^{-n+2}\bigr)\,\prod_{k>j}\bigl(x_k-t^{-n+1}\bigr)}
{e^*_{n-1}(x;t)},
\quad
e^*_{n-1}(x;t)=\sum_{j=1}^n\prod_{k<j}\bigl(x_k-t^{-n+2}\bigr)\prod_{k>j}\bigl(x_k-t^{-n+1}\bigr).
\]

\smallskip
\noindent\textbf{(Step 1) Push-TASEP move.} The particle at site $j$, say with
label $a$, is activated and travels clockwise by the $t$-Push TASEP rule: if
there are $m$ weaker particles in the system (labels $<a$, vacancies counting as
label $0$), the active particle moves to the $k$-th such weaker particle with
probability $t^{k-1}/[m]_t$. If that site held a positive label, that particle
becomes active and repeats; the cascade ends when the active particle reaches a
vacancy. Afterwards site $j$ is vacant; we treat this vacancy as a new particle
labelled $a:=0$.

\smallskip
\noindent\textbf{(Step 2) Vacancy sweep.} The label-$0$ particle goes to site $1$
and travels clockwise. Reaching a site $k$ ($1\le k\le j-1$) holding a label
$b\ge0$, it skips that site with probability $1-p_k$ if $b\ge a$ and $1-q_k$ if
$b<a$; otherwise it settles there, activating/displacing the resident, which then
continues clockwise toward site $j$ by the same rule. The active particle stops
once it displaces another or arrives at site $j$, where it settles.
\end{qthreedefinition}

We denote the transition probabilities of Steps 1 and 2 by $P^{(1)}_j$ and
$P^{(2)}_j$, so that for $\mu,\nu\in\qthreeSnl$,
\[
P(\mu,\nu)=\sum_{1\le j\le n}P_j\sum_{\rho\in\qthreeSnl:\rho_j=0}
P^{(1)}_j(\mu,\rho)\,P^{(2)}_j(\rho,\nu).
\]

\begin{qthreeremark}[nontriviality]
The rates $P_j,p_k,q_k$ are explicit rational functions of $x_1,\dots,x_n$ and
$t$ only. They do \emph{not} reference $F^*_\mu$, so the chain is nontrivial in
the required sense (unlike a Metropolis--Hastings chain built from the target
distribution).
\end{qthreeremark}

\begin{qthreetheorem}\label{thm:main}
In the interpolation $t$-Push TASEP with content $\lambda$, the stationary
probability of $\mu\in\qthreeSnl$ is $\pi^*_\lambda(\mu)=F^*_\mu(x;1,t)/P^*_\lambda(x;1,t)$.
\end{qthreetheorem}

\section*{Two-line queues and the proof}

We use classical two-line queues \cite{CMW22} and signed two-line queues
\cite{BDW25a}, specialized to $q=1$. Let $\mathcal Q^\eta_\kappa$ be the classical
two-line queues with top row $\eta$ and bottom row $\kappa$, with pair-weight
generating function $a^\eta_\kappa=\sum_{Q\in\mathcal Q^\eta_\kappa}\qthreewt_{\mathrm{pair}}(Q)$.
Let $\mathcal G^\alpha_\mu$ be the signed two-line queues with top row $\alpha\in\qthreeZ^n$
and bottom row $\mu$, with $b^\alpha_\mu=\sum_{Q\in\mathcal G^\alpha_\mu}\qthreewt_{\mathrm{pair}}(Q)$,
and full weight $\qthreewt(Q)=\qthreewt_{\mathrm{pair}}(Q)\qthreewt_{\mathrm{ball}}(Q)$ where
\begin{equation}\label{eq:wta}
\qthreewt_\alpha\,b^\alpha_\mu=\sum_{Q\in\mathcal G^\alpha_\mu}\qthreewt(Q),
\quad \qthreewt_\alpha:=\prod_{k:\alpha_k>0}x_k\prod_{k:\alpha_k<0}\bigl(-t^{-(n-1)}\bigr).
\end{equation}

\begin{qthreedefinition}\label{def:unsigned}
For a signed two-line queue $Q\in\mathcal G^\alpha_\mu$ let $\overline Q$ be its
unsigned version (forget signs on the top row); its top composition is
$\|\alpha\|=(|\alpha_1|,\dots,|\alpha_n|)$. Let $\mathcal G^\kappa_\mu$ be the set of
unsigned objects so obtained, with $\|\alpha\|=\kappa$. For $\overline Q\in\mathcal G^\kappa_\mu$
assign a nontrivial pairing $p$ the weight $(1-t)t^{\mathrm{skip}(p)}$, and a top-row
ball $B$ in column $k$ labelled $a>0$ with the ball below labelled $b$ the weight
\[
\qthreewt(B)=\begin{cases}x_k-t^{-(n-1)} & b=a,\\ x_k & b>a,\\ t^{-(n-1)} & b<a.\end{cases}
\]
\end{qthreedefinition}

\begin{qthreelemma}\label{lem:signforget}
For $\lambda$ with distinct parts and $\kappa,\mu\in\qthreeSnl$, and $\overline Q\in\mathcal G^\kappa_\mu$,
$\qthreewt(\overline Q)=\sum_Q\qthreewt(Q)$, the sum over signed $Q$ forgetting to $\overline Q$.
\end{qthreelemma}

\begin{proof}
Enumerate the admissible sign choices for each top-row ball $B$ labelled $a>0$:
a $-$ sign is forced when the ball below is a vacancy or labelled $b<a$; a $+$
sign is forced when $b>a$; both are allowed when $b=a$. One checks this matches
\eqref{eq:wta}: a $-$ sign multiplies the ball weight by $-1$ in passing from
$Q$ to $\overline Q$, but the weight of the pairing attached to $B$ is also
multiplied by $-1$, so signs cancel consistently.
\end{proof}

For $\kappa\in\qthreeSnl$ set $c^\kappa_\mu:=\sum_{\alpha:\|\alpha\|=\kappa}\qthreewt_\alpha\,b^\alpha_\mu$.
Combining \eqref{eq:wta} and Lemma~\ref{lem:signforget}:

\begin{qthreelemma}\label{lem:ckm}
$c^\kappa_\mu=\sum_{\overline Q\in\mathcal G^\kappa_\mu}\qthreewt(\overline Q)$. Since $\lambda$ has distinct
parts, $\mathcal G^\kappa_\mu$ is empty or a single element.
\end{qthreelemma}

\subsection*{The closing recursion (where the hypotheses on $\lambda$ enter)}

Fix a weakly order-preserving $\phi:\qthreeN\to\qthreeN$ and partitions $\lambda,\kappa$ with
$\phi(\lambda)=\kappa$. For $\eta\in S_n(\kappa)$ set
$G^*_\eta(x;t):=\sum_{\rho\in\qthreeSnl:\,\phi(\rho)=\eta}F^*_\rho(x;1,t)$. The following
is an interpolation analogue of a result of Ayyer--Martin--Williams
\cite{AMW25}, provable in the same way using \cite{AS19}.

\begin{qthreetheorem}\label{thm:G}
For $\lambda,\kappa$ as above and all $\eta\in S_n(\kappa)$, at $q=1$,
\[
\frac{G^*_\eta(x;t)}{P^*_\lambda(x;1,t)}=\frac{F^*_\eta(x;1,t)}{P^*_\kappa(x;1,t)} .
\]
\end{qthreetheorem}

\begin{qthreecorollary}\label{cor:rec}
Let $\rho$ be a composition with $\rho_i\ne1$ for all $i$, with $k$ nonzero
parts, and set $\eta=\rho^-$ (where $\rho^-_i=\max(\rho_i-1,0)$). Then at $q=1$,
\[
F^*_\rho(x;1,t)=F^*_\eta(x;1,t)\cdot e^*_k(x;t).
\]
\end{qthreecorollary}

\begin{proof}
Let $\lambda,\kappa$ reorder $\rho,\eta$. Take $\phi:i\mapsto\max(i-1,0)$, so
$\phi(\rho)=\eta$. \emph{Because $\lambda$ has no part of size $1$}, $\phi$ is a
bijection from $\{0,2,3,\dots\}$ onto $\{0,1,2,\dots\}$, so $\rho$ is the unique
composition in $\qthreeSnl$ with $\phi(\rho)=\eta$; hence $G^*_\eta=F^*_\rho$. By
Theorem~\ref{thm:G},
$F^*_\rho/P^*_\lambda=F^*_\eta/P^*_\kappa$. Finally, by \eqref{eq:factor} and the
fact that $\kappa$ is obtained from $\lambda$ by deleting the largest column (of
size $k$), $P^*_\lambda/P^*_\kappa=e^*_k(x;t)$, giving the claim.
\end{proof}

This is exactly where the restriction on $\lambda$ is essential: \emph{no part of
size $1$} makes $\phi$ bijective so the recursion closes, and the \emph{unique
part of size $0$} guarantees $\eta$ has a unique $0$-part so the column-deletion
is unambiguous. The hypotheses are structural, not (as one might guess) a device
to avoid a $0/0$ normalization.

\subsection*{From two-line queues to the dynamics}

\begin{qthreeproposition}\label{prop:step2}
For $\rho,\nu\in\qthreeSnl$ with $\rho_j=0$,
\[
P^{(2)}_j(\rho,\nu)=\frac{c^\rho_\nu}{\prod_{k<j}(x_k-t^{-n+2})\prod_{k>j}(x_k-t^{-n+1})},
\quad\text{equivalently}\quad
P_j\,P^{(2)}_j(\rho,\nu)=\frac{c^\rho_\nu}{e^*_{n-1}} .
\]
\end{qthreeproposition}

\begin{proof}[Idea]
Step 2 is encoded by the unique element of $\mathcal G^\rho_\nu$ (Lemma~\ref{lem:ckm}):
a particle labelled $a$ moving from column $k$ to $k'$ corresponds to a
nontrivial pairing joining a top ball in column $k$ to a bottom ball in column
$k'$; stationary particles give trivial pairings. Dividing each ball/pairing
weight in $\qthreewt(\overline Q)$ by the matching factor of
$D=\prod_{k<j}(x_k-t^{-n+2})\prod_{k>j}(x_k-t^{-n+1})$ reproduces, case by case,
the skip/settle probabilities $1-p_k,\,1-q_k,\,p_k,\,q_k$ of Step 2 (and a
factor $1$ for $k>j$, which never moves). Hence $\qthreewt(\overline Q)/D=P^{(2)}_j(\rho,\nu)$.
\end{proof}

\begin{qthreeproposition}\label{prop:transition}
If $\lambda$ is restricted and $\mu,\nu\in\qthreeSnl$, then
$\displaystyle P(\mu,\nu)=\sum_{\rho\in\qthreeSnl}\frac{a^\mu_\rho\,c^\rho_\nu}{e^*_{n-1}}$.
\end{qthreeproposition}

\begin{proof}
Combine the encoding of Step 1 by classical two-line queues
\cite[Lem.~5.4]{AMW25} (giving $P^{(1)}_j(\mu,\rho)$ via $a^\mu_\rho$) with
Proposition~\ref{prop:step2}:
\[
P(\mu,\nu)=\sum_j P_j\!\!\sum_{\rho:\rho_j=0}\!P^{(1)}_j(\mu,\rho)P^{(2)}_j(\rho,\nu)
=\sum_j\sum_{\rho:\rho_j=0}\frac{a^\mu_\rho c^\rho_\nu}{e^*_{n-1}}
=\sum_{\rho\in\qthreeSnl}\frac{a^\mu_\rho c^\rho_\nu}{e^*_{n-1}}.
\]
\end{proof}

\subsection*{Proof of Theorem~\ref{thm:main}}

\begin{proof}
Fix $\nu\in\qthreeSnl$. By \cite[Thm.~1.15, Lem.~5.6]{BDW25a},
\[
F^*_\nu(x;1,t)=\sum_{\eta\in\qthreeN^n}F^{*\eta}_\nu(x;t)\,F^*_{\eta^-}(x;1,t),
\qquad
F^{*\eta}_\nu(x;t)=\sum_{\kappa\in\qthreeN^n}a^\eta_\kappa\,c^\kappa_\nu .
\]
By Corollary~\ref{cor:rec} (using that $\eta$ has a unique $0$-part),
$F^*_{\eta^-}(x;1,t)=F^*_\eta(x;1,t)/e^*_{n-1}(x;t)$. Substituting,
\[
F^*_\nu(x;1,t)=\sum_{\eta\in\qthreeN^n}F^*_\eta(x;1,t)\sum_{\kappa\in\qthreeN^n}
\frac{a^\eta_\kappa\,c^\kappa_\nu}{e^*_{n-1}(x;t)}
=\sum_{\eta\in\qthreeN^n}F^*_\eta(x;1,t)\,P(\eta,\nu),
\]
the last equality by Proposition~\ref{prop:transition}. Thus the vector
$\bigl(F^*_\mu(x;1,t)\bigr)_\mu$ is left-fixed by the transition matrix, i.e.\
proportional to the stationary distribution. Since
$P^*_\lambda=\sum_{\mu\in\qthreeSnl}F^*_\mu$, normalizing gives
$\pi^*_\lambda(\mu)=F^*_\mu(x;1,t)/P^*_\lambda(x;1,t)$.
\end{proof}

\section*{Remark on the failure mode to avoid}

Replacing this chain by the ordinary multispecies ASEP (matrix-product ansatz /
Zamolodchikov--Faddeev algebra) does \emph{not} solve the stated problem: that
construction yields the ordinary Macdonald polynomials as stationary weights
\cite{CMW22,AMW25}, not the \emph{interpolation} polynomials $F^*_\mu,P^*_\lambda$.
The interpolation correction is precisely what the bell $P_j$ and the
push-with-sweep dynamics supply, and the proof above runs through signed
two-line queues rather than a trace over an auxiliary algebra.

\begin{thebibliography}{9}
\bibitem{AS19} P.~Alexandersson, M.~Sawhney.
  \emph{Properties of non-symmetric Macdonald polynomials at $q=1$ and $q=0$}.
  Ann.\ Comb.\ \textbf{23} (2019), 219--239.
\bibitem{AMW25} A.~Ayyer, J.~Martin, L.~Williams.
  \emph{The inhomogeneous $t$-PushTASEP and Macdonald polynomials at $q=1$}.
  Ann.\ Inst.\ Henri Poincar\'e D, 2025.
\bibitem{BDW25a} H.~Ben Dali, L.~Williams.
  \emph{A combinatorial formula for interpolation Macdonald polynomials}.
  Preprint, arXiv:2510.02587, 2025.
\bibitem{BDW26} H.~Ben Dali, L.~Williams.
  \emph{A probabilistic interpretation for interpolation Macdonald polynomials}.
  Preprint, arXiv:2602.13492, 2026.
\bibitem{CMW22} S.~Corteel, O.~Mandelshtam, L.~Williams.
  \emph{From multiline queues to Macdonald polynomials via the exclusion process}.
  Amer.\ J.\ Math.\ \textbf{144} (2022), 395--436.
\bibitem{Dolega} M.~Dolega.
  \emph{Strong factorization property of Macdonald polynomials and higher-order
  Macdonald's positivity conjecture}.
  J.\ Algebraic Combin.\ \textbf{46} (2017), 135--163.
\end{thebibliography}

\endgroup
% ===== End q3_solution.tex =====

% ===== Begin q4_solution.tex =====
\clearpage
\section*{Problem 4}
\addcontentsline{toc}{section}{Problem 4}
\begingroup
\hypersetup{colorlinks=true,
            linkcolor=blue!50!black,
            citecolor=blue!50!black,
            urlcolor=blue!50!black}
\newcommand{\qfourR}{\mathbb{R}}
\newcommand{\qfourZ}{\mathbb{Z}}
\newcommand{\qfourbn}{\boxplus_n}
\newcommand{\qfourScore}{\mathsf{S}^{\mathrm{ff}}_n}
\newcommand{\qfourInfo}{\mathcal{K}^{\mathrm{ff}}_n}
\newcommand{\qfourRootMap}{\Psi_{\qfourbn}}
\newcommand{\qfourJac}{\mathsf{D}_{\qfourbn}}
\newcommand{\qfourVar}{\operatorname{Var}}
\newcommand{\qfourHess}{\operatorname{Hess}}
\newcommand{\qfourPSD}{\succeq 0}
\newcommand{\qfoursgn}{\operatorname{sgn}}
\newtheorem{qfourtheorem}{Theorem}
\newtheorem{qfourlemma}{Lemma}
\newtheorem{qfourproposition}{Proposition}
\newtheorem{qfourcorollary}{Corollary}
\theoremstyle{definition}
\newtheorem{qfourobservation}{Observation}
\title{Solution to Problem 4: The Finite Free Stam Inequality}
\author{}
\date{}
\maketitle


\section*{Statement and Answer}

For a monic real-rooted polynomial $p$ of degree $n$ with roots
$\alpha=(\alpha_1,\dots,\alpha_n)$, define the \emph{finite free score vector}
$\qfourScore(\alpha)\in(\qfourR\cup\{\infty\})^n$ by
\[
\qfourScore(\alpha)_i \;=\; \sum_{j\ne i}\frac{1}{\alpha_i-\alpha_j},
\]
and the \emph{finite free Fisher information}
$\qfourInfo(p):=\|\qfourScore(\alpha)\|^2$, with $\qfourInfo(p)=\infty$ if $p$ has a
multiple root. Let $\qfourbn$ denote the finite free convolution of degree-$n$
polynomials.

\medskip
\noindent\textbf{Theorem.} \emph{For any two monic real-rooted polynomials
$p,q$ of degree $n$,}
\[
\frac{1}{\qfourInfo(p\qfourbn q)}\;\ge\;\frac{1}{\qfourInfo(p)}+\frac{1}{\qfourInfo(q)}.
\]

\medskip
\noindent The answer is \textbf{YES}. The result was conjectured by
D.~Shlyakhtenko; the proof below follows Blachman's strategy for the classical
Stam inequality, with the key analytic input replaced by a convexity statement
for hyperbolic polynomials.

\section*{Preliminaries on $\qfourbn$}

We recall the standard description of the finite free convolution. If $\alpha$
and $\beta$ are vectors of roots for $p$ and $q$, then
\begin{equation}\label{eq:walsh}
p(x)\qfourbn q(x)\;=\;\sum_{\pi\in S_n}\prod_{i=1}^n\bigl(x-\alpha_i-\beta_{\pi(i)}\bigr),
\end{equation}
and Walsh's theorem guarantees that $p\qfourbn q$ is again real-rooted
\cite{Walsh, MSS22}. Thus $\qfourbn$ induces a map
\[
\qfourRootMap:\qfourR^n\times\qfourR^n\to\qfourR^n,
\]
where $\qfourRootMap(\alpha,\beta)$ is the ordered root vector of $p\qfourbn q$. Write
$m_k(\alpha)=\frac1n\sum_i\alpha_i^k$ and
$\qfourVar(\alpha)=m_2(\alpha)-m_1(\alpha)^2$, and let $\mathbf 1_n$ be the all-ones
vector.

\begin{qfourproposition}[properties of $\qfourbn$]\label{prop:props}
If $\gamma=\qfourRootMap(\alpha,\beta)$ then
\begin{enumerate}\setlength\itemsep{2pt}
\item (additivity) $m_1(\gamma)=m_1(\alpha)+m_1(\beta)$ and
$\qfourVar(\gamma)=\qfourVar(\alpha)+\qfourVar(\beta)$;
\item (translation) $\qfourRootMap(\alpha+t\mathbf 1_n,\beta)=\gamma+t\mathbf 1_n$
and likewise in the second argument.
\end{enumerate}
\end{qfourproposition}

\begin{proof}
(i) follows from \eqref{eq:walsh} together with the Newton identities applied
to the first two coefficients of $p\qfourbn q$; (ii) is immediate from
\eqref{eq:walsh}.
\end{proof}

\section*{Score vectors via the heat flow}

The link between $\qfourInfo$ and the dynamics of roots is the following, due to
Shlyakhtenko.

\begin{qfourlemma}[score vector as a velocity]\label{lem:heat}
Assume $p$ has simple roots. Let $p_t(x):=\exp\!\bigl(-\tfrac{t}{2}\partial_x^2\bigr)p(x)$
and let $\alpha(t)$ be the ordered root vector of $p_t$. Then
$\alpha_i'(0)=\sum_{j\ne i}\frac{1}{\alpha_i-\alpha_j}$; that is,
$\alpha'(0)=\qfourScore(\alpha)$.
\end{qfourlemma}

\begin{proof}
Roots stay simple near $t=0$. Implicitly differentiating
$p(\alpha_i(t))-\tfrac{t}{2}p''(\alpha_i(t))+t^2R(\alpha_i(t),t)=0$ at $t=0$
gives $\alpha_i'(0)=\tfrac12\,p''(\alpha_i)/p'(\alpha_i)$, which equals
$\sum_{j\ne i}(\alpha_i-\alpha_j)^{-1}$.
\end{proof}

The essential point is that $\exp(-\tfrac t2\partial_x^2)$ \emph{commutes with}
$\qfourbn$ (finite free convolution commutes with any constant-coefficient
differential operator \cite{MSS22}). We exploit this twice.

\section*{Reduction to simple roots}

\begin{qfourlemma}\label{lem:simple}
Let $T_\epsilon:=(1-\epsilon\,d/dx)^n$. For monic real-rooted $p$ of degree $n$:
(i) $T_\epsilon p$ is monic real-rooted with simple roots; (ii)
$\qfourInfo(T_\epsilon p)\to\qfourInfo(p)$ as $\epsilon\to0$; (iii)
$(T_\epsilon p)\qfourbn(T_\epsilon q)=T_\epsilon^2(p\qfourbn q)$.
\end{qfourlemma}

\begin{proof}
(i) is Nuij's theorem \cite{Nuij}; (ii) holds since $\qfourInfo$ is continuous in
the roots and the roots are continuous in $\epsilon$; (iii) holds because $\qfourbn$
commutes with the differential operator $1-\epsilon\,d/dx$ \cite{MSS22}.
\end{proof}

By Lemma~\ref{lem:simple}(iii) and (ii), it suffices to prove the Theorem when
$p$, $q$, and $p\qfourbn q$ all have simple roots; we assume this henceforth. Then
$\alpha,\beta,\gamma=\qfourRootMap(\alpha,\beta)$ have distinct entries and are
smooth functions of the coefficients of the corresponding polynomials. Let
$\qfourJac$ denote the Jacobian of $\qfourRootMap$ at $(\alpha,\beta)$.

\section*{Step 1: score vectors and the Jacobian}

\begin{qfourobservation}\label{obs:scores}
For all $a,b\ge0$,
\[
\qfourJac\bigl(a\,\qfourScore(\alpha),\,b\,\qfourScore(\beta)\bigr)\;=\;(a+b)\,\qfourScore(\gamma).
\]
\end{qfourobservation}

\begin{proof}
Set $p_t=\exp(-\tfrac t2\partial_x^2)p$ with ordered roots $\alpha(t)$, and
define $q_t,\beta(t)$ and $r_t,\gamma(t)$ analogously. Because $\qfourbn$ commutes
with the heat operator,
\[
r_{(a+b)t}\;=\;p_{at}\qfourbn q_{bt},
\]
so $\gamma((a+b)t)=\qfourRootMap(\alpha(at),\beta(bt))$. Differentiating at
$t=0$ and using the chain rule,
\[
(a+b)\,\gamma'(0)\;=\;\qfourJac\!\begin{pmatrix} a\,\alpha'(0)\\ b\,\beta'(0)\end{pmatrix}.
\]
Lemma~\ref{lem:heat} identifies $\alpha'(0)=\qfourScore(\alpha)$, etc., giving the
claim.
\end{proof}

\section*{Step 2: the Jacobian is a contraction on mean-zero vectors}

This is the substance of the proof.

\begin{qfourproposition}\label{prop:contraction}
If $u,v\in\qfourR^n$ are each orthogonal to $\mathbf 1_n$, then
\[
\|\qfourJac(u,v)\|^2\;\le\;\|u\|^2+\|v\|^2.
\]
\end{qfourproposition}

Write $\gamma_i=(\qfourRootMap)_i(\alpha,\beta)$ for the coordinate functions and
$H^{(i)}:=\qfourHess(\qfourRootMap)_i$. Let
$V=\{(u,v)\in\qfourR^n\times\qfourR^n:u^\ast\mathbf 1_n=v^\ast\mathbf 1_n=0\}$.

\begin{qfourlemma}[Hessian identity]\label{lem:hess}
For all $w\in V$,
\[
w^\ast \qfourJac\qfourJac^\ast w
\;=\;w^\ast\Bigl(I_n\oplus I_n-\sum_{i=1}^n\gamma_i H^{(i)}\Bigr)w.
\]
\end{qfourlemma}

\begin{proof}
Fix $w=(u,v)\in V$ and set $\alpha(t)=\alpha+tu$, $\beta(t)=\beta+tv$,
$\gamma(t)=\qfourRootMap(\alpha(t),\beta(t))$. By the additivity of variance
(Proposition~\ref{prop:props}(i)),
\[
m_2(\gamma(t))-m_1(\gamma(t))^2
=m_2(\alpha(t))+m_2(\beta(t))-m_1(\alpha(t))^2-m_1(\beta(t))^2 .
\]
Since $(u,v)\in V$, the means $m_1(\alpha(t)),m_1(\beta(t))$ — hence
$m_1(\gamma(t))$ — are constant in $t$. Differentiating twice at $t=0$,
\[
\partial_t^2 m_2(\gamma(t))\big|_0=\partial_t^2\bigl(m_2(\alpha(t))+m_2(\beta(t))\bigr)\big|_0 .
\]
The left side equals
$\tfrac1n\sum_i\partial_t^2(\gamma_i^2)
=\tfrac2n\sum_i\bigl((\partial_t\gamma_i)^2+\gamma_i\partial_t^2\gamma_i\bigr)
=\tfrac2n\bigl(w^\ast \qfourJac\qfourJac^\ast w+\sum_i\gamma_i w^\ast H^{(i)}w\bigr)$,
while the right side equals $\tfrac2n\,w^\ast w$. Rearranging gives the
identity.
\end{proof}

The minus sign means we must show $\sum_i\gamma_i H^{(i)}\qfourPSD$ on $V$. This is
where real-rootedness of $\qfourbn$ is used in an essential way, via the theory of
hyperbolic polynomials.

\begin{qfourtheorem}[Bauschke--G\"uler--Lewis--Sendov \cite{BGLS}]\label{thm:bgls}
Let $f\in\qfourR[x_1,\dots,x_m]$ be hyperbolic in the direction $w\in\qfourR^m$, and for
$a\in\qfourR^m$ let $\lambda_1(a)\ge\dots\ge\lambda_m(a)$ be the roots of
$t\mapsto f(a+tw)$. Then for each $k$, $\sigma_k(a):=\sum_{i=1}^k\lambda_i(a)$ is
convex in $a$.
\end{qfourtheorem}

\begin{qfourcorollary}\label{cor:psd}
For any reals $c_1\ge\dots\ge c_n$, the matrix $\sum_{i=1}^n c_i H^{(i)}$ is PSD.
\end{qfourcorollary}

\begin{proof}
Consider $f(x,a_1,\dots,a_n,b_1,\dots,b_n)=\sum_{\pi\in S_n}\prod_{i=1}^n
(x-a_i-b_{\pi(i)})$. By \eqref{eq:walsh} and Walsh's theorem $f$ is homogeneous
and hyperbolic in the direction $e_1=(1,0,\dots,0)$. By Theorem~\ref{thm:bgls}
each $\sigma_k(x,a,b)=\sum_{i=1}^k\lambda_i(x,a,b)$ is convex, where
$\lambda_i(x,a,b)$ are the roots of $f((x,a,b)+te_1)$. Since the $c_i$ are
ordered,
\[
L(x,a,b):=\sum_{i=1}^n c_i\lambda_i(x,a,b)
\]
is a nonnegative combination of the $\sigma_k$ (by Abel summation), hence convex;
so $\qfourHess L=\sum_i c_i\qfourHess\lambda_i\qfourPSD$ everywhere. Evaluating at $x=0$,
$a=\alpha$, $b=\beta$, we have $\qfourHess\lambda_i=H^{(i)}$, giving
$\sum_i c_i H^{(i)}\qfourPSD$.
\end{proof}

\begin{proof}[Proof of Proposition~\ref{prop:contraction}]
Let $(u,v)\in V$. By Lemma~\ref{lem:hess},
\[
\|\qfourJac(u,v)\|^2=(u,v)^\ast \qfourJac\qfourJac^\ast(u,v)
=\|u\|^2+\|v\|^2-\sum_{i=1}^n\gamma_i\,(u,v)^\ast H^{(i)}(u,v).
\]
The roots $\gamma_1\ge\dots\ge\gamma_n$ are ordered, so
Corollary~\ref{cor:psd} (with $c_i=\gamma_i$) gives $\sum_i\gamma_i H^{(i)}\qfourPSD$,
whence $\sum_i\gamma_i(u,v)^\ast H^{(i)}(u,v)\ge0$. The claim follows.
\end{proof}

\section*{Step 3: Blachman's argument}

\begin{proof}[Proof of the Theorem]
First, $\qfourScore(\alpha)\perp\mathbf 1_n$: indeed
$\sum_i\sum_{j\ne i}(\alpha_i-\alpha_j)^{-1}=0$ since each unordered pair
contributes $+$ and $-$ once. Likewise $\qfourScore(\beta)\perp\mathbf 1_n$. So
Proposition~\ref{prop:contraction} applies to
$(a\,\qfourScore(\alpha),\,b\,\qfourScore(\beta))$:
\[
\bigl\|\qfourJac\bigl(a\,\qfourScore(\alpha),b\,\qfourScore(\beta)\bigr)\bigr\|^2
\;\le\;a^2\|\qfourScore(\alpha)\|^2+b^2\|\qfourScore(\beta)\|^2 .
\]
Combining with Observation~\ref{obs:scores},
\[
(a+b)^2\,\qfourInfo(p\qfourbn q)\;\le\;a^2\,\qfourInfo(p)+b^2\,\qfourInfo(q),
\]
valid for all $a,b\ge0$. Choosing $a=1/\qfourInfo(p)$ and $b=1/\qfourInfo(q)$, the right
side becomes $1/\qfourInfo(p)+1/\qfourInfo(q)$ and the factor $(a+b)^2$ becomes
$\bigl(1/\qfourInfo(p)+1/\qfourInfo(q)\bigr)^2$, so
\[
\Bigl(\tfrac{1}{\qfourInfo(p)}+\tfrac{1}{\qfourInfo(q)}\Bigr)^2\qfourInfo(p\qfourbn q)
\;\le\;\tfrac{1}{\qfourInfo(p)}+\tfrac{1}{\qfourInfo(q)} .
\]
Dividing through yields
$\dfrac{1}{\qfourInfo(p\qfourbn q)}\ge\dfrac{1}{\qfourInfo(p)}+\dfrac{1}{\qfourInfo(q)}$.
The simple-root reduction (Lemma~\ref{lem:simple}) extends this to all
real-rooted $p,q$.
\end{proof}

\section*{Remarks}

The proof mirrors Blachman's \cite{Blachman} derivation of the classical Stam
inequality \cite{Stam}: there the score of a sum is a conditional expectation of
the marginal score and the inequality reduces to Cauchy--Schwarz. Here the role
of that conditional-expectation step is played by
Observation~\ref{obs:scores} (the score of $p\qfourbn q$ is the image of the marginal
scores under the Jacobian $\qfourJac$) together with the operator-norm bound
Proposition~\ref{prop:contraction}. Real-rootedness of $\qfourbn$ is used exactly
once and essentially — through the hyperbolicity of \eqref{eq:walsh} and the
BGLS convexity theorem — which is why the inequality fails for convolutions that
do not preserve real-rootedness.

\begin{thebibliography}{9}
\bibitem{MSS22} A.~W.~Marcus, D.~A.~Spielman, N.~Srivastava.
  \emph{Finite free convolutions of polynomials}.
  Probab.\ Theory Related Fields \textbf{182} (2022), 807--848.
\bibitem{Walsh} J.~L.~Walsh.
  \emph{On the location of the roots of certain types of polynomials}.
  Trans.\ Amer.\ Math.\ Soc.\ \textbf{24} (1922), 163--180.
\bibitem{Nuij} W.~Nuij.
  \emph{A note on hyperbolic polynomials}.
  Math.\ Scand.\ \textbf{23} (1968), 69--72.
\bibitem{Blachman} N.~Blachman.
  \emph{The convolution inequality for entropy powers}.
  IEEE Trans.\ Inform.\ Theory \textbf{11} (1965), 267--271.
\bibitem{BGLS} H.~H.~Bauschke, O.~G\"uler, A.~S.~Lewis, H.~S.~Sendov.
  \emph{Hyperbolic polynomials and convex analysis}.
  Canad.\ J.\ Math.\ \textbf{53} (2001), 470--488.
\bibitem{Stam} A.~Stam.
  \emph{Some inequalities satisfied by the quantities of information of Fisher
  and Shannon}. Inform.\ Control \textbf{2} (1959), 101--112.
\end{thebibliography}

\endgroup
% ===== End q4_solution.tex =====

% ===== Begin q5_solution.tex =====
\clearpage
\section*{Problem 5}
\addcontentsline{toc}{section}{Problem 5}
\begingroup
\hypersetup{colorlinks=true,
            linkcolor=blue!50!black,
            citecolor=blue!50!black,
            urlcolor=blue!50!black}
\newcommand{\qfiveaO}{\mathcal{O}}
\newcommand{\qfiveSp}{\mathrm{Sp}}
\newcommand{\qfiveZ}{\mathbb{Z}}
\newcommand{\qfiveR}{\mathbb{R}}
\newcommand{\qfivegconn}{\operatorname{gconn}}
\newcommand{\qfiveconn}{\operatorname{conn}}
\newcommand{\qfiveSub}{\operatorname{Sub}}
\newcommand{\qfiveMap}{\operatorname{Map}}
\newtheorem{qfivetheorem}{Theorem}
\newtheorem{qfivelemma}{Lemma}
\newtheorem{qfiveproposition}{Proposition}
\newtheorem{qfivecorollary}{Corollary}
\theoremstyle{definition}
\newtheorem{qfivedefinition}{Definition}
\newtheorem{qfiveremark}{Remark}
\title{Solution to Problem 5: The $\qfiveaO$-Slice Filtration for $N_\infty$ Operads}
\author{}
\date{}
\maketitle


\section*{Setup: transfer systems and indexing systems}

Fix a finite group $G$. Following Balchin--Barnes--Roitzheim \cite{BBR} and
Rubin \cite{Rubin}, the homotopy theory of $N_\infty$ operads on $G$ is governed
by the combinatorics of \emph{transfer systems}.

\begin{qfivedefinition}[transfer system]
A \emph{transfer system} on $G$ is a partial order, denoted $\to$, on the set
$\qfiveSub(G)$ of subgroups of $G$ satisfying:
\begin{enumerate}\setlength\itemsep{2pt}
\item it refines subgroup inclusion: if $H\to K$ then $H\subseteq K$;
\item it is conjugation invariant: if $H\to K$ and $g\in G$ then
$gHg^{-1}\to gKg^{-1}$;
\item it is closed under restriction: if $H\to K$ and $J\subseteq K$ then
$H\cap J\to J$.
\end{enumerate}
Transfer systems on $G$ form a poset under refinement.
\end{qfivedefinition}

\begin{qfivedefinition}[admissible sets, indexing system]
Let $\qfiveaO$ be a transfer system on $G$. A finite $H$-set $T=\coprod_i H/K_i$ is
\emph{admissible} for $\qfiveaO$ if $K_i\to H$ for all $i$. The collection $\qfiveaO(H)$ of
admissible $H$-sets, as $H$ varies, is the associated \emph{indexing system}; we
also denote the corresponding $N_\infty$ operad by $\qfiveaO$.
\end{qfivedefinition}

For a finite $G$-set $T$ we have the \emph{$T$-norm}
$N^T:\qfiveSp^G\to\qfiveSp^G$, built from the Hill--Hopkins--Ravenel norms \cite{HHR16}
by $N^{G/H}(E)=N^G_H i_H^\ast(E)$ and $N^{T_0\sqcup T_1}(E)=N^{T_0}(E)\otimes
N^{T_1}(E)$.

\section*{The $\qfiveaO$-slice filtration}

We use equivariant localization (nullification) to construct the relative slice
tower. Recall that the equivariant localizing subcategory generated by a set of
objects is its closure under homotopy colimits, retracts, and tensors with orbit
spectra.

\begin{qfivedefinition}[$\qfiveaO$-slice connective spectra]
For an indexing system $\qfiveaO$, let $\tau^\qfiveaO_{\ge n}$ be the equivariant
localizing subcategory of $\qfiveSp^G$ generated by
\[
\bigl\{\, G_+\otimes_H N^T S^1 \;\bigm|\; T\in\qfiveaO(H),\ |T|\ge n \,\bigr\}.
\]
This is the category of \emph{$\qfiveaO$-slice $n$-connective} spectra.
\end{qfivedefinition}

\begin{qfiveremark}
There is an equivariant homeomorphism $N^T S^1\cong S^{\qfiveR\cdot T}$, the
representation sphere of the permutation representation of $T$. So the
$n$-connective spectra are generated by representation spheres of admissible
sets of cardinality $\ge n$. This replaces the regular-representation slice cells
$G_+\wedge_H S^{m\rho_H}$ of the complete (Hill--Hopkins--Ravenel) filtration:
in the incomplete setting one is only permitted the norms indexed by sets that
$\qfiveaO$ actually admits.
\end{qfiveremark}

\begin{qfivedefinition}
For an indexing system $\qfiveaO$ and $n\ge 0$:
\begin{itemize}\setlength\itemsep{2pt}
\item the $n$th \emph{$\qfiveaO$-slice truncation} $P^n_\qfiveaO:\qfiveSp^G_{\ge 0}\to
\qfiveSp^G_{\ge 0}$ is the nullification killing $\tau^\qfiveaO_{\ge(n+1)}$;
\item the $n$th \emph{$\qfiveaO$-slice cover} $P^\qfiveaO_n$ is the fiber of
$\mathrm{Id}\Rightarrow P^{n-1}_\qfiveaO$;
\item the $n$th \emph{$\qfiveaO$-slice} $P^n_{n,\qfiveaO}(E)$ is the fiber of
$P^n_\qfiveaO(E)\to P^{n-1}_\qfiveaO(E)$.
\end{itemize}
A $G$-spectrum $E$ is an \emph{$\qfiveaO$-$n$-slice} if $E\in\tau^\qfiveaO_{\ge n}$ and
$E\to P^n_\qfiveaO E$ is an equivalence.
\end{qfivedefinition}

That the tower converges (for $E$ connective, $E\xrightarrow{\sim}\varprojlim
P^n_\qfiveaO E$) follows because the Postnikov connectivity of $G_+\otimes_H N^T S^1$
is $|T/H|\to\infty$ with $|T|$. The suspension $\Sigma N^T S^1\simeq N^{T\sqcup
\ast}S^1$ shows $\Sigma:\tau^\qfiveaO_{\ge k}\to\tau^\qfiveaO_{\ge(k+1)}$, whence the
$\infty$-category of $\qfiveaO$-$k$-slices is discrete.

\section*{Geometric fixed points and the characteristic function}

Stable equivalences in $\qfiveSp^G$ are detected on geometric fixed points
$\Phi^H$ for all $H\subseteq G$, so the connectivity of geometric fixed points
is the governing invariant.

\begin{qfivedefinition}[geometric connectivity]
For $E\in\qfiveSp^G$, let $\qfivegconn(E):\qfiveSub(G)\to\qfiveZ\cup\{\pm\infty\}$ be
$\qfivegconn(E)(H):=\qfiveconn(\Phi^H E)$.
\end{qfivedefinition}

\begin{qfivedefinition}[characteristic function]
For a transfer system $\qfiveaO$, define $\chi^\qfiveaO:\qfiveSub(G)\to\qfiveSub(G)$ by
\[
\chi^\qfiveaO(H)\;:=\;\min\{K\mid K\to H\}\;=\;\bigcap_{K\to H}K .
\]
\end{qfivedefinition}

The index $[G:\chi^\qfiveaO(G)]$ is the maximal cardinality of an admissible $G$-set
with a single orbit; equivalently $[G:\chi^\qfiveaO(H)]$ measures how large an
admissible $H$-set $\qfiveaO$ permits per orbit. This subgroup, \emph{not} $|H|$
itself, is what enters the connectivity bound.

\section*{Main Theorem}

\begin{qfivetheorem}\label{thm:main}
Let $\qfiveaO$ be a transfer system on $G$ and let $E$ be a connective $G$-spectrum.
Then $E\in\tau^\qfiveaO_{\ge n}$ if and only if for all $H\subseteq G$,
\begin{equation}\label{eq:crit}
[H:\chi^\qfiveaO(H)]\cdot\qfivegconn(E)(H)\;\ge\;n .
\end{equation}
\end{qfivetheorem}

Equivalently, writing
$n=\min_{H\subseteq G}\,[H:\chi^\qfiveaO(H)]\cdot\qfivegconn(E)(H)$, one has
$E\in\tau^\qfiveaO_{\ge n}$; this $n$ is the \emph{$\qfiveaO$-slice connectivity} of $E$.

When $\qfiveaO$ is complete, $\chi^\qfiveaO(H)=\{e\}$ for all $H$, $[H:\chi^\qfiveaO(H)]=|H|$,
and \eqref{eq:crit} recovers the Hill--Yarnall reformulation \cite{HY18} of the
Hill--Hopkins--Ravenel slice connectivity.

\section*{Proof}

\subsection*{Forward direction (necessity)}

\begin{qfivelemma}\label{lem:forward}
If $E\in\tau^\qfiveaO_{\ge n}$ then $[H:\chi^\qfiveaO(H)]\cdot\qfivegconn(E)(H)\ge n$ for all
$H\subseteq G$.
\end{qfivelemma}

\begin{proof}
By restriction it suffices to treat $H=G$. Geometric fixed points preserve
homotopy colimits and extensions and annihilate induced spectra from proper
subgroups, so it suffices to check the generators $N^T S^1$ for admissible
$G$-sets $T$ with $|T|\ge n$. Decompose $T=\sum_H n_H\, G/H$. Then
$\Phi^G(N^T S^1)=S^{|T/G|}$ with $|T/G|=\sum_H n_H$, while
\[
|T|=\sum_H n_H[G:H]\ge n .
\]
By definition $[G:\chi^\qfiveaO(G)]$ is the maximal $[G:H]$ with $G/H\in\qfiveaO(\ast)$
(and all others divide it), so
\[
[G:\chi^\qfiveaO(G)]\cdot\sum_H n_H \;\ge\; \sum_H n_H[G:H]\;\ge\; n,
\]
i.e.\ $[G:\chi^\qfiveaO(G)]\cdot\qfivegconn(N^T S^1)(G)\ge n$, as required. When
$\chi^\qfiveaO(G)=\{e\}$ this is exactly \cite[Thm.~2.5]{HY18}.
\end{proof}

\subsection*{Converse via isotropy separation}

Let $\mathcal P$ be the family of proper subgroups of $G$ and consider the
isotropy separation sequence
\[
E\mathcal P_+\otimes E\;\longrightarrow\; E\;\longrightarrow\;
\widetilde{E\mathcal P}\otimes E .
\]
We argue by downward induction on the subgroup lattice. The first piece is built
from inductions $G/H_+\otimes E$ for $H\in\mathcal P$, controlled by the
inductive hypothesis on restrictions; the second piece is \emph{geometric} and
is handled by the following.

\begin{qfivelemma}[geometric slices]\label{lem:geom}
Let $M$ be a geometric Mackey functor. For any $\qfiveaO$,
$\Sigma^k HM$ is a $k\cdot[G:\chi^\qfiveaO(G)]$-$\qfiveaO$-slice.
\end{qfivelemma}

\begin{proof}
Since $M$ is geometric, for any finite $G$-set $T$ the inclusion of fixed points
$S^{|T/G|}\hookrightarrow N^T S^1$ induces an equivalence
$S^{|T/G|}\otimes HM\xrightarrow{\sim} N^T S^1\otimes HM$. For the lower bound,
choose an admissible $T$ with $|T|$ maximal subject to $|T/G|=k$, namely
$T=k\,G/\chi^\qfiveaO(G)$ (as $\chi^\qfiveaO(G)$ is the minimal subgroup with
$\chi^\qfiveaO(G)\to G$); this gives $\Sigma^k HM\in\tau^\qfiveaO_{\ge k[G:\chi^\qfiveaO(G)]}$.
For the upper bound, if $T$ is admissible with $|T|>k[G:\chi^\qfiveaO(G)]$ then
$|T/G|>k$, so
\[
[N^T S^1,\Sigma^k HM]^G\cong[\Phi^G N^T S^1,\Sigma^k HM(G/G)]
\cong[S^{|T/G|},\Sigma^k HM(G/G)]=0 .
\]
Hence $\Sigma^k HM$ is exactly a $k[G:\chi^\qfiveaO(G)]$-slice.
\end{proof}

\begin{qfivelemma}\label{lem:induced}
Let $\mathcal F$ be a family and $\tau$ an equivariant localizing subcategory. If
$i_H^\ast E\in i_H^\ast\tau$ for all $H\in\mathcal F$, then $E\mathcal
F_+\otimes E\in\tau$.
\end{qfivelemma}

\begin{proof}
$E\mathcal F_+\otimes E$ lies in the localizing subcategory generated by
$G/H_+\otimes E\cong G_+\otimes_H i_H^\ast E$ for $H\in\mathcal F$; by hypothesis
each such object is in $\tau$. (Same argument as \cite[Lem.~2.4]{HY18}.)
\end{proof}

\begin{proof}[Proof of Theorem~\ref{thm:main}, converse]
Suppose $E$ is connective with $[H:\chi^\qfiveaO(H)]\cdot\qfivegconn(E)(H)\ge n$ for all
$H$. In the isotropy separation sequence, Lemma~\ref{lem:geom} (applied to the
homotopy Mackey functors of the geometric part $\widetilde{E\mathcal P}\otimes
E$) shows its $\qfiveaO$-slice connectivity is $\ge n$. By downward induction on the
subgroup lattice, the hypothesis on proper subgroups together with
Lemma~\ref{lem:induced} shows $E\mathcal P_+\otimes E$ also has $\qfiveaO$-slice
connectivity $\ge n$. Since localizing subcategories are closed under
extensions, $E\in\tau^\qfiveaO_{\ge n}$.
\end{proof}

This is precisely the $\qfiveaO$-relative form of \cite[Thm.~2.5]{HY18}; the only
new ingredient is that the multiplier $|H|$ is replaced by $[H:\chi^\qfiveaO(H)]$,
reflecting that $\qfiveaO$ admits norms only up to index $[H:\chi^\qfiveaO(H)]$ per orbit.

\section*{Examples and consistency checks}

\paragraph{Complete $\qfiveaO$.} Here every $G/H$ is admissible, $\chi^\qfiveaO(H)=\{e\}$,
and \eqref{eq:crit} reads $|H|\cdot\qfivegconn(E)(H)\ge n$ for all $H$ --- the
Hill--Hopkins--Ravenel slice connectivity in the Hill--Yarnall formulation
\cite{HY18, HHR16}.

\paragraph{Trivial $\qfiveaO$.} No nontrivial transfers: $\chi^\qfiveaO(H)=H$ for all $H$,
so $[H:\chi^\qfiveaO(H)]=1$ and \eqref{eq:crit} becomes $\qfivegconn(E)(H)\ge n$ for all
$H$. The only admissible sets are trivial orbits, and the filtration is the
geometric Postnikov tower.

\paragraph{$G=\qfiveZ/p$.} The two transfer systems are trivial and complete. For the
complete one, $\chi^\qfiveaO(\qfiveZ/p)=\{e\}$ gives the bound $p\cdot\qfivegconn(E)(\qfiveZ/p)\ge n$
at the top together with $\qfivegconn(E)(e)\ge n$; for the trivial one both indices
are $1$ and the conditions are $\qfivegconn(E)(H)\ge n$. These agree with the
$C_p$-slice computations of \cite{HY18}.

\section*{Scope}

The model-categorical and $\infty$-categorical foundations for the
norm functors and equivariant localization are taken from \cite{HHR16, BBR,
Rubin}; the geometric-Mackey-functor input (Lemma~\ref{lem:geom}) adapts
\cite[\S6]{HillPrimer}. The argument is the Hill--Yarnall template
\cite{HY18} carried out with the indexing-system-restricted norms in place of
the full regular-representation cells, and the characteristic function
$\chi^\qfiveaO$ tracking the available norms.

\begin{thebibliography}{9}
\bibitem{HHR16} M.~A.~Hill, M.~J.~Hopkins, D.~C.~Ravenel.
  \emph{On the nonexistence of elements of Kervaire invariant one}.
  Ann.\ of Math.\ (2) \textbf{184} (2016), 1--262.
\bibitem{HY18} M.~A.~Hill, C.~Yarnall.
  \emph{A new formulation of the equivariant slice filtration with applications
  to $C_p$-slices}. Proc.\ Amer.\ Math.\ Soc.\ \textbf{146} (2018), 3605--3614.
\bibitem{HillPrimer} M.~A.~Hill.
  \emph{The equivariant slice filtration: a primer}.
  Homology Homotopy Appl.\ \textbf{14} (2012), 143--166.
\bibitem{BBR} S.~Balchin, D.~Barnes, C.~Roitzheim.
  \emph{$N_\infty$-operads and associahedra}.
  Pacific J.\ Math.\ \textbf{315} (2021), 285--304.
\bibitem{Rubin} J.~Rubin.
  \emph{Detecting Steiner and linear isometries operads}.
  Glasg.\ Math.\ J.\ \textbf{63} (2021), 307--342.
\bibitem{Wilson} D.~Wilson.
  \emph{On categories of slices}. Preprint, arXiv:1711.03472, 2017.
\end{thebibliography}

\endgroup
% ===== End q5_solution.tex =====

% ===== Begin q6_solution.tex =====
\clearpage
\section*{Problem 6}
\addcontentsline{toc}{section}{Problem 6}
\begingroup
\hypersetup{colorlinks=true,
            linkcolor=blue!50!black,
            citecolor=blue!50!black,
            urlcolor=blue!50!black}
\newcommand{\qsixR}{\mathbb{R}}
\newcommand{\qsixE}{\mathbb{E}}
\newcommand{\qsixPr}{\mathbb{P}}
\newcommand{\qsixeps}{\varepsilon}
\newcommand{\qsixTr}{\operatorname{Tr}}
\newcommand{\qsixTrs}{\operatorname{Tr}_{\sigma}}
\newcommand{\qsixdiag}{\operatorname{diag}}
\newcommand{\qsixopnorm}[1]{\left\lVert #1 \right\rVert}
\newcommand{\qsixdel}{\delta}
\newcommand{\qsixSpan}{\operatorname{span}}
\newcommand{\qsixim}{\operatorname{im}}
\newcommand{\qsixrank}{\operatorname{rank}}
\newcommand{\qsixtL}{\widetilde{L}}
\newcommand{\qsixLt}{L^{\dagger}}
\newcommand{\qsixLth}{L^{\dagger/2}}
\newcommand{\qsixBarrier}[2]{\mathcal{B}^{#1}_{#2}}
\newtheorem{qsixtheorem}{Theorem}
\newtheorem{qsixlemma}[qsixtheorem]{Lemma}
\newtheorem{qsixproposition}[qsixtheorem]{Proposition}
\newtheorem{qsixclaim}[qsixtheorem]{Claim}
\title{$\qsixeps$-Light Subsets in Graphs}
\author{}
\date{}
\maketitle


\section*{Statement and Answer}

Let $G=(V,E,w)$ be a weighted graph on $n=|V|$ vertices with Laplacian
$L=\sum_{(s,t)\in E}w(s,t)\,(\delta_s-\delta_t)(\delta_s-\delta_t)^{\!\top}$.
For $S\subseteq V$ let $G_S=(V,E(S,S))$ keep only edges with both endpoints in
$S$ and let $L_S$ be its Laplacian (same ambient dimension as $L$). $S$ is
\emph{$\qsixeps$-light} iff $\qsixeps L - L_S\succeq 0$.

\medskip
\noindent\textbf{Theorem.} \emph{For every weighted graph $G=(V,E,w)$ and every
$\qsixeps\in(0,1)$, there exists an $\qsixeps$-light subset $S\subseteq V$ with
$|S|\ge \lfloor \qsixeps n/42\rfloor$.}

\medskip
\noindent The answer is \textbf{YES}, with the explicit constant $c=1/42$.

\section*{Notation and basic facts}

Write $\qsixLt$ for the Moore--Penrose pseudoinverse and $\qsixLth$ for its PSD square
root. Define
\[
  \qsixtL_S      := \qsixLth L_S \qsixLth, \qquad
  \qsixtL_{S,T}  := \qsixLth L_{S,T}  \qsixLth, \qquad
  \qsixtL_{S,t}  := \qsixLth L_{S,\{t\}}\qsixLth.
\]
Then $\qsixLth L \qsixLth = \Pi$, the orthogonal projector onto $\qsixSpan\{L\}$, and the
condition $\qsixeps L\succeq L_S$ is equivalent (by conjugation with $\qsixLth$ on
$\qsixim L$) to the operator-norm bound
\begin{equation}\label{eq:opnorm-form}
  \qsixopnorm{\qsixtL_S}_{\mathrm{op}} \;\le\; \qsixeps.
\end{equation}

\medskip
\noindent\emph{Leverage scores.}\quad The leverage score of an edge $(s,t)$ is
\[
  \ell(s,t)
   := w(s,t)\,(\delta_s-\delta_t)^{\!\top}\qsixLt(\delta_s-\delta_t)
    = \qsixTr\!\bigl(L_{\{s\},\{t\}}\qsixLt\bigr),
\]
and zero on non-edges. For $S,T\subseteq V$ set
$\ell(S,T):=\sum_{s\in S,t\in T}\ell(s,t)$ and write
$\ell(S):=\ell(S,V\setminus S)$ for the boundary leverage. By linearity
$\ell(S,T)=\qsixTr(\qsixtL_{S,T})$. Foster's identity, in the form
$\sum_{(s,t)\in E}\ell(s,t)=n-1$ for connected $G$, follows from
$\sum_{(s,t)}L_{\{s\},\{t\}}=L$ and $\qsixTr(L\qsixLt)=\qsixrank L = n-1$.

\medskip
\noindent\emph{Bounded-rank upper barrier.}\quad For a symmetric $A$ with
eigenvalues $\lambda_1\ge\cdots\ge\lambda_n$ and a real $u$, define
\[
  \qsixBarrier{u}{\sigma}(A) \;:=\;
  \begin{cases}
    \displaystyle\sum_{i=1}^{\sigma}\dfrac{1}{u-\lambda_i} & \text{if } u>\lambda_1,\\[2pt]
    \infty & \text{otherwise,}
  \end{cases}
\]
and the top-$\sigma$ trace $\qsixTrs(A):=\sum_{i=1}^{\sigma}\lambda_i$. We write
$\qsixBarrier{u}{\sigma}(S)$ for $\qsixBarrier{u}{\sigma}(\qsixtL_S)$, and note
$\qsixBarrier{u}{\sigma}(A)=\qsixTrs\bigl((uI-A)^{-1}\bigr)$ when $u>\lambda_1$.

For the rest of this note we fix
\[
  \qsixdel := \frac{21}{n}, \qquad
  \phi := \frac{n}{21}, \qquad
  \sigma := \lfloor \qsixeps n/42 \rfloor.
\]

\section*{Proof}

We build $S$ greedily, adding one vertex at a time and increasing the barrier
parameter $u$ by $\qsixdel$ at each step. Throughout, two invariants are
maintained: the bounded-rank barrier stays at most $\phi$, and the boundary
leverage stays at most $4|S|$.

\subsection*{Step 1: Initial set and target}

Pick any $v_0\in V$ and set $S_0:=\{v_0\}$, $u_0:=\qsixeps/2$. Since $G_{S_0}$ has
no edges, $\qsixtL_{S_0}=0$, and
\[
  \qsixBarrier{u_0}{\sigma}(S_0) \;=\; \frac{\sigma}{u_0}
   \;\le\; \frac{\qsixeps n/42}{\qsixeps/2} \;=\; \frac{n}{21} \;=\; \phi.
\]
After $\sigma$ greedy steps, the barrier will sit at $u_\sigma=u_0+\sigma\qsixdel$.
We chose parameters so that
\[
  u_\sigma = \tfrac{\qsixeps}{2} + \sigma\cdot\tfrac{21}{n}
           \le \tfrac{\qsixeps}{2} + \tfrac{\qsixeps}{2} = \qsixeps,
\]
so once we reach $|S|=\sigma+1$ with $\qsixBarrier{u_\sigma}{\sigma}(S)$ finite, every
eigenvalue of $\qsixtL_S$ lies below $\qsixeps$ and \eqref{eq:opnorm-form} gives
$\qsixeps L\succeq L_S$.

\subsection*{Step 2: Greedy lemma}

\begin{qsixlemma}\label{lem:greedy}
Suppose $|S|\le\sigma$, $\ell(S)\le 4|S|$, and $\qsixBarrier{u}{\sigma}(S)\le\phi$. Then
there exists $t\notin S$ with
\[
  \qsixBarrier{u+\qsixdel}{\sigma}(S\cup\{t\})\le\phi
  \qquad\text{and}\qquad
  \ell(S\cup\{t\})\le\ell(S)+4.
\]
\end{qsixlemma}

\noindent The lemma is proved by two averaging arguments: more than half the
candidates $t$ keep the leverage budget (Step 3), and at least half keep the
barrier (Step 4). The two sets therefore intersect.

\subsection*{Step 3: Leverage-score averaging}

\begin{qsixclaim}\label{cl:Tlev}
For every $T\subseteq V$, $\sum_{t\in T}\ell(t,T)\le 2(|T|-1)$.
\end{qsixclaim}
\begin{proof}
By symmetry, $\sum_{t\in T}\ell(t,T)=2\qsixTr(L_T\qsixLt)$. Since $L_T\preceq L$, all
eigenvalues of $L_T\qsixLt$ lie in $[0,1]$, and $L_T$ has rank at most $|T|-1$, so
$\qsixTr(L_T\qsixLt)\le|T|-1$.
\end{proof}

For $t\notin S$,
\[
  \ell(S\cup\{t\}) \le \ell(S\cup\{t\},V\setminus S)
                   = \ell(S)+\ell(t,V\setminus S),
\]
so it suffices to control $\ell(t,V\setminus S)$. Claim~\ref{cl:Tlev} applied
to $T=V\setminus S$ gives
$\sum_{t\notin S}\ell(t,V\setminus S)<2|V\setminus S|$, whence by Markov more
than half of $t\notin S$ satisfy $\ell(t,V\setminus S)\le 4$.

\subsection*{Step 4: Barrier averaging}

This is the technical heart. The barrier-update lemma (Lemma~\ref{lem:bu}
below) reduces our task to showing that for at least half of $t\notin S$,
\begin{equation}\label{eq:U}
  U(S,t) \;:=\;
  \frac{\qsixTrs(M^{-2}\qsixtL_{S,t})}{\qsixBarrier{u}{\sigma}(S)-\qsixBarrier{u+\qsixdel}{\sigma}(S)}
   + \qsixTrs(M^{-1}\qsixtL_{S,t}) \;<\; 1,
\end{equation}
where $M:=(u+\qsixdel)I-\qsixtL_S$. We claim
\begin{equation}\label{eq:total}
  \sum_{t\notin S} U(S,t) \;\le\; \tfrac{5}{\qsixdel} + 5\phi.
\end{equation}
Granting \eqref{eq:total}, since each $U(S,t)\ge 0$, more than half of
$t\notin S$ satisfy
\[
  U(S,t)\le\frac{2}{n-|S|}\!\left(\tfrac{5}{\qsixdel}+5\phi\right)
        \le \frac{2}{41n/42}\cdot\tfrac{10n}{21} < 1
\]
for $|S|\le\sigma\le n/42$.

\medskip
\noindent\emph{Proof of \eqref{eq:total}.}\quad Let $\Pi_S$ project onto
$\qsixim\qsixtL_S$, and $\Pi_S^\perp:=I-\Pi_S$. Since
$M=(u+\qsixdel)(\Pi_S+\Pi_S^\perp)-\qsixtL_S$ and $\Pi_S^\perp\qsixtL_S=0$,
\[
  M^p \;=\; (u+\qsixdel)^p\Pi_S^\perp + \bigl((u+\qsixdel)\Pi_S-\qsixtL_S\bigr)^p,
  \qquad p\in\{-1,-2\}.
\]
By Ky Fan subadditivity of $\qsixTrs$ \cite[Thm.~4.3.47a]{HJ12},
\begin{equation}\label{eq:split}
  \qsixTrs(M^p\qsixtL_{S,t}) \;\le\;
  \qsixTrs\!\bigl((u+\qsixdel)^p\Pi_S^\perp\qsixtL_{S,t}\bigr)
  + \qsixTrs\!\Bigl(((u+\qsixdel)\Pi_S-\qsixtL_S)^p\qsixtL_{S,t}\Bigr).
\end{equation}

\smallskip
\emph{$\Pi_S^\perp$ piece.}\quad
$\qsixTrs(\Pi_S^\perp\qsixtL_{S,t})\le\qsixTr(\qsixtL_{S,t})=\ell(S,t)$ (the PSD product has
non-negative eigenvalues), and summing,
\[
  \sum_{t\notin S}\qsixTrs\!\bigl((u+\qsixdel)^p\Pi_S^\perp\qsixtL_{S,t}\bigr)
   \le (u+\qsixdel)^p\,\ell(S)\le (u+\qsixdel)^p\cdot 4|S|.
\]

\smallskip
\emph{$\Pi_S$ piece.}\quad
$\sum_{t\notin S}\qsixtL_{S,t}=\qsixtL_{S,V\setminus S}\preceq\Pi$, and
$((u+\qsixdel)\Pi_S-\qsixtL_S)^p$ is PSD on $\qsixim\Pi_S$, so
\[
\begin{aligned}
\sum_{t\notin S}\qsixTrs\!\Bigl(((u+\qsixdel)\Pi_S-\qsixtL_S)^p\qsixtL_{S,t}\Bigr)
  &\le \qsixTr\!\Bigl(((u+\qsixdel)\Pi_S-\qsixtL_S)^p\qsixtL_{S,V\setminus S}\Bigr)\\
  &\le \qsixTr\!\Bigl(((u+\qsixdel)\Pi_S-\qsixtL_S)^p\Bigr)
   = \qsixTr(\Pi_S M^p\Pi_S) \le \qsixTrs(M^p),
\end{aligned}
\]
the last step by the Ky Fan max principle.

\smallskip
Since all eigenvalues of $M$ are at most $u+\qsixdel$, for $p<0$ all eigenvalues of
$M^p$ are at least $(u+\qsixdel)^p$, so $\qsixTrs(M^p)\ge\sigma(u+\qsixdel)^p\ge|S|(u+\qsixdel)^p$.
Combining,
\[
  \sum_{t\notin S}\qsixTrs(M^p\qsixtL_{S,t}) \;\le\; 4|S|(u+\qsixdel)^p + \qsixTrs(M^p)
   \;\le\; 5\,\qsixTrs(M^p).
\]
Now we use (Claim~\ref{cl:phi-diff})
$\qsixBarrier{u}{\sigma}(S)-\qsixBarrier{u+\qsixdel}{\sigma}(S)\ge\qsixdel\,\qsixTrs(M^{-2})$ to bound the
first term of $U(S,t)$, and $\qsixTrs(M^{-1})=\qsixBarrier{u}{\sigma}(S)\le\phi$ for the
second:
\[
  \sum_{t\notin S}U(S,t)
   \;\le\; \frac{5\,\qsixTrs(M^{-2})}{\qsixdel\,\qsixTrs(M^{-2})} + 5\,\qsixBarrier{u}{\sigma}(S)
   \;\le\; \tfrac{5}{\qsixdel}+5\phi,
\]
which is \eqref{eq:total}.\hfill$\square$

\subsection*{Step 5: Combining and conclusion}

Iterating Lemma~\ref{lem:greedy} $\sigma$ times starting from $S_0$ yields
$S$ with $|S|=\sigma+1$, $\ell(S)\le 4\sigma$, and
$\qsixBarrier{u_\sigma}{\sigma}(S)<\infty$, hence (by Step~1) every eigenvalue of
$\qsixtL_S$ is at most $u_\sigma\le\qsixeps$, so $\qsixeps L\succeq L_S$. This proves the
theorem with $|S|\ge\lfloor\qsixeps n/42\rfloor$.\qed

\section*{Auxiliary lemmas}

\begin{qsixlemma}[Barrier update]\label{lem:bu}
Let $A,B\succeq 0$ and $\qsixdel>0$. Set $M:=(u+\qsixdel)I-A$. If
$\qsixBarrier{u}{\sigma}(A)<\infty$ and
\[
  \frac{\qsixTrs(M^{-2}B)}{\qsixBarrier{u}{\sigma}(A)-\qsixBarrier{u+\qsixdel}{\sigma}(A)}
   + \qsixTrs(M^{-1}B) \;<\; 1,
\]
then $\qsixBarrier{u+\qsixdel}{\sigma}(A+B)\le\qsixBarrier{u}{\sigma}(A)$.
\end{qsixlemma}

\begin{proof}
$\qsixBarrier{u}{\sigma}(A)<\infty$ gives $u>\lambda_{\max}(A)$, hence
$M,M^{-1},M^{-2}\succ 0$. Write $B=CC^{\!\top}$. By Sherman--Morrison--Woodbury,
\[
  (M-CC^{\!\top})^{-1}
   = M^{-1} + M^{-1}C(I-C^{\!\top}M^{-1}C)^{-1}C^{\!\top}M^{-1}.
\]
Since $\qsixTrs(M^{-1}B)=\qsixTrs(C^{\!\top}M^{-1}C)<1$ the inverse exists and the
second summand is PSD, so $\lambda_{\max}(A+B)<u+\qsixdel$. By Ky Fan
subadditivity \cite[Thm.~4.3.47a]{HJ12},
\[
  \qsixBarrier{u+\qsixdel}{\sigma}(A+B)
   \le \qsixTrs(M^{-1})
       + \qsixTrs\!\bigl((I-C^{\!\top}M^{-1}C)^{-1}C^{\!\top}M^{-2}C\bigr).
\]
Since $\qsixopnorm{C^{\!\top}M^{-1}C}\le\qsixTrs(C^{\!\top}M^{-1}C)<1$, the operator
$(I-C^{\!\top}M^{-1}C)^{-1}\preceq(1-\qsixTrs(C^{\!\top}M^{-1}C))^{-1}I$, and
monotonicity of $\qsixTrs$ on PSD-ordered products gives
\[
  \qsixTrs\!\bigl((I-C^{\!\top}M^{-1}C)^{-1}C^{\!\top}M^{-2}C\bigr)
   \le \frac{\qsixTrs(M^{-2}B)}{1-\qsixTrs(M^{-1}B)}.
\]
Combining with
$\qsixTrs(M^{-1})=\qsixBarrier{u}{\sigma}(A)-(\qsixBarrier{u}{\sigma}(A)-\qsixBarrier{u+\qsixdel}{\sigma}(A))$
and the hypothesis gives the claimed bound.
\end{proof}

\begin{qsixclaim}\label{cl:phi-diff}
$\qsixBarrier{u}{\sigma}(S)-\qsixBarrier{u+\qsixdel}{\sigma}(S)\ge\qsixdel\,\qsixTrs(M^{-2})$.
\end{qsixclaim}
\begin{proof}
For each top-$\sigma$ eigenvalue $\lambda_i$ of $\qsixtL_S$,
\[
  \tfrac{1}{u-\lambda_i}-\tfrac{1}{u+\qsixdel-\lambda_i}
   = \tfrac{\qsixdel}{(u-\lambda_i)(u+\qsixdel-\lambda_i)}
   \ge \tfrac{\qsixdel}{(u+\qsixdel-\lambda_i)^2};
\]
sum over $i$.
\end{proof}

\section*{Remarks}

\begin{enumerate}\setlength\itemsep{4pt}
\item \textbf{Constant.} The constant $1/42$ comes from the parameter choices
$\qsixdel=21/n$, $\phi=n/21$, which balance the $5/\qsixdel$ and $5\phi$ contributions
in \eqref{eq:total}. It is not believed optimal; improving it is open.
\item \textbf{Key innovation.} The standard upper-barrier function tracks all
eigenvalues; here we track only the top $\sigma$. This adaptation, together
with the Ky Fan inequalities for $\qsixTrs$, is what makes the argument go
through.
\item \textbf{Connection to spectral sparsification.} An $\qsixeps$-light $S$ with
$|S|\ge c\qsixeps n$ says the induced subgraph on $S$ contributes at most $\qsixeps L$
spectrally, giving a perspective on ``light'' as ``spectrally independent''.
\item \textbf{Vertex vs.\ edge sampling.} The edge-partitioning analogue
admits an optimal constant of $1/2$ via the Kadison--Singer resolution. The
vertex-light variant is more delicate because edges sharing a vertex are
correlated; the deterministic barrier method bypasses concentration entirely.
\end{enumerate}

\begin{thebibliography}{9}
\bibitem{HJ12} R.~A.~Horn and C.~R.~Johnson.
  \emph{Matrix Analysis}, 2nd ed.\ Cambridge University Press, 2012.
\end{thebibliography}

\endgroup
% ===== End q6_solution.tex =====

% ===== Begin q7_solution.tex =====
\clearpage
\section*{Problem 7}
\addcontentsline{toc}{section}{Problem 7}
\begingroup
\hypersetup{colorlinks=true,
            linkcolor=blue!50!black,
            citecolor=blue!50!black,
            urlcolor=blue!50!black}
\newcommand{\qsevenR}{\mathbb{R}}
\newcommand{\qsevenQ}{\mathbb{Q}}
\newcommand{\qsevenZ}{\mathbb{Z}}
\newcommand{\qsevenLgp}{\mathbb{L}}
\newcommand{\qsevensign}{\operatorname{sign}}
\newcommand{\qsevencd}{\operatorname{cd}}
\newcommand{\qseventM}{\widetilde M}
\newcommand{\qsevenim}{\operatorname{im}}
\newtheorem{qseventheorem}{Theorem}
\newtheorem{qsevenlemma}{Lemma}
\newtheorem{qsevenproposition}{Proposition}
\title{Uniform Lattices with $2$-Torsion as Fundamental Groups}
\author{}
\date{}
\maketitle


\section*{Statement and Answer}

Let $\Gamma$ be a uniform (cocompact) lattice in a real semisimple Lie group
$G$, and suppose $\Gamma$ contains an element of order $2$.

\medskip
\noindent\textbf{Theorem.} \emph{Such a $\Gamma$ cannot be the fundamental group
of a closed manifold without boundary whose universal cover is acyclic over
$\qsevenQ$.}

\medskip
\noindent The answer is \textbf{NO}.

\medskip
A torsion-free uniform lattice in $G$ \emph{is} the fundamental group of such a
manifold --- a closed locally symmetric space, whose universal cover is the
contractible symmetric space $X=G/K$. The content of the theorem is that
$2$-torsion obstructs this. Fowler, in his Ph.D.\ thesis, proved the
corresponding statement when $\Gamma$ has \emph{odd} torsion (a proof appears in
\cite{W2}); the present argument extends Fowler's theorem to the case of
$2$-torsion, following Cappell--Weinberger--Yan.

We stress at the outset what does \emph{not} work: a proof using only that $M$
is a finite Poincar\'e complex over $\qsevenQ$ cannot succeed, because such complexes
with these fundamental groups exist (Proposition~\ref{prop:counterexample}). The
argument must use that $M$ is a genuine manifold, and it does so through the
symmetric signature and an equivariant fundamental-class computation.

\section*{What a finite-complex / Poincar\'e-duality argument cannot see}

It is tempting to argue purely with rational Poincar\'e duality and Euler
characteristics. That route is blocked by an explicit construction of Fowler.

\begin{qsevenproposition}\label{prop:counterexample}
There exist groups of the form $\pi_1(M_3)\times\Gamma$, with $\Gamma$ a lattice
with torsion, that are fundamental groups of spaces having the rational homotopy
type of a finite complex satisfying rational Poincar\'e duality.
\end{qsevenproposition}

\begin{proof}[Construction]
Let $\Gamma$ be a lattice in a linear semisimple group $G$, and choose a
homomorphism $\Gamma\to\Delta$ to a finite group with torsion-free kernel
$\Gamma_0$. Let $E\Delta$ be a contractible free $\Delta$-space and $M_3$ any
closed hyperbolic $3$-manifold. The space
\[
  M_3\times\bigl(K\backslash G/\Gamma_0\times E\Delta\bigr)/\Delta
\]
has the rational homotopy type of a finite complex and satisfies rational
Poincar\'e duality; its fundamental group is $\pi_1(M_3)\times\Gamma$, a lattice
in $\mathrm{SO}(3,1)\times G$.
\end{proof}

In particular, ``multiplicativity of the Euler characteristic in finite covers''
\emph{fails} for infinite complexes with finitely generated rational homology: a
rationally acyclic example with rational Euler characteristic $1$ is the
universal cover of $\qsevenR P^2\vee(\bigvee_\infty S^2)$ after equivariantly attaching
cells $D^3\times\qsevenZ/2$ along a free $\qsevenZ[1/2][\qsevenZ/2]$-basis of
$\pi_2=\qsevenZ[-1]\oplus\qsevenZ[\qsevenZ/2]^{\infty}$ to kill the homology; both it and its
$\qsevenZ/2$-quotient become rationally acyclic with rational Euler characteristic $1$.
Hence the obstruction cannot be an Euler-characteristic denominator, and any
proof relying on $\chi(M)\in\qsevenZ$ versus $\chi_\qsevenQ(\Gamma)\notin\qsevenZ$ is invalid. We
also note that any proof resting on a naive Lefschetz fixed-point argument is
unsound: the relevant Lefschetz-type statement is \emph{false} (witness
$\qsevenR^1$ with $f$ a translation, which has no fixed points yet Lefschetz number
$-1$ in that sense).

\section*{Proof of the Theorem}

\subsection*{Reduction}

Passing to a normal subgroup of finite index, it suffices to treat
$\Gamma=\pi\rtimes\qsevenZ_2$ for a torsion-free $\pi$, a lattice in $G$, such that
there is an involution on $M_0:=K\backslash G/\pi$ --- by isometries of the
locally symmetric metric --- whose fixed set $F$ is nonempty. (If $F$ is
disconnected one argues componentwise; we write the connected case.) Here $M_0$
is a closed aspherical manifold, a $K(\pi,1)$.

Suppose, for contradiction, that $X^m$ is a closed manifold with
$\pi_1(X)\cong\Gamma$ and rationally acyclic universal cover. Let $Y\to X$ be the
$2$-fold cover corresponding to $\pi\le\Gamma$; then $\pi_1(Y)\cong\pi$ and the
universal cover of $Y$ (the same as that of $X$) is rationally acyclic.

\subsection*{Step 1: a degree-one fundamental class}

Work with symmetric signatures in the symmetric $L$-group $L(\qsevenR\pi)$ over the
reals. (Over $\qsevenR$ the symmetric and quadratic $L$-groups agree, so we may write
$L(\qsevenR\pi)$ unambiguously.) There is a map $f:Y\to M_0$ inducing the identity on
$\pi_1=\pi$. It need not have degree one, but over $\qsevenR$ every degree has a square
root, so the symmetric signature is sensitive only to the sign of the degree,
and $f$ induces an equivalence of symmetric signatures.

Because the Novikov conjecture holds for $\pi$ (a lattice in a semisimple group;
it satisfies the Borel/Novikov package via its proper action on the
nonpositively curved $X$), the assembly map
\[
  H_m\bigl(B\pi;\qsevenLgp(\qsevenR)\bigr)\;\longrightarrow\;L_m(\qsevenR\pi)
\]
is injective. The symmetric signature of $Y$ is detected in the degree-$m$ piece
$H_m(B\pi;\qsevenZ)$ as the image of the fundamental class. Injectivity forces
\begin{equation}\label{eq:deg-one}
  f_*[Y]\;=\;[M_0]\quad\text{in}\quad H_m(B\pi;\qsevenQ),
\end{equation}
i.e.\ $f$ is rationally degree one onto the aspherical $M_0$.

\subsection*{Step 2: an equivariant cobordism and the fixed set}

We now run the cobordism argument of \cite{W1}. Consider any closed manifold $Z$
with fundamental group $\pi$ carrying an involution inducing the given
automorphism of $\pi$, together with the image of $[Z]$ in $H_m(B\Gamma;\qsevenZ_2)$.
Standard equivariant homotopy theory gives an equivariant map $g:Z\to M_0$,
hence a map of fixed sets $Z^{\qsevenZ_2}\to F$. We claim
\begin{equation}\label{eq:fixed}
  g_*[Z]\;=\;g_*[Z^{\qsevenZ_2}],
\end{equation}
using the map $\qsevenZ_2\times\pi_1F\to\Gamma$ and the periodicity in the group
homology of $\qsevenZ_2$ to raise the dimension from $\dim F$ to $\dim M_0$.

The relevant cobordism realizing \eqref{eq:fixed} runs between $Z$ and a
projective-space bundle over $Z^{\qsevenZ_2}$ --- the projectivized normal bundle of
$Z^{\qsevenZ_2}$ --- whose fundamental class is the asserted element by the
Leray--Hirsch theorem. Explicitly it is $Z\times[0,1]$, with $Z\times\{1\}$
modified by quotienting the free $\qsevenZ/2$-action on the complement of an
equivariant regular neighbourhood of $Z^{\qsevenZ_2}$.

\subsection*{Step 3: the contradiction}

Apply Step 2 in two ways.

\emph{For $Y$.} The deck $\qsevenZ_2$-action on $Y$ (equivalently the action with
quotient $X$) is \emph{free}, so $Y^{\qsevenZ_2}=\varnothing$ and the right-hand side
of \eqref{eq:fixed} is $0$. Hence the corresponding image of the fundamental
class is $0$.

\emph{For $M_0$.} The isometric $\qsevenZ_2$-action on $M_0$ has fixed set $F$ which is
aspherical, and the Borel construction shows that
$\qsevenZ_2\times F\to\Gamma$ induces an injection on homology in dimension
$\dim(M_0/\qsevenZ_2)$ (an isomorphism in higher dimensions; see \cite{Borel}). Since
the fundamental class of an aspherical manifold is nontrivial in its group
homology, the corresponding image for $M_0$ is \emph{nonzero}.

But \eqref{eq:deg-one} identifies the $Y$-image with the $M_0$-image rationally:
the degree-one map $f$ carries the (zero) class for $Y$ to the (nonzero) class
for $M_0$. This is a contradiction. Therefore no such closed manifold $X$
exists, and $\Gamma$ is not the fundamental group of a closed manifold with
rationally acyclic universal cover.\qed

\section*{Remarks}

\begin{enumerate}\setlength\itemsep{4pt}
\item \textbf{Where the manifold hypothesis enters.} The argument uses that $X$
is a manifold twice: to form the symmetric signature and apply assembly
injectivity (Step 1), and to run the equivariant normal-bundle cobordism (Step
2). Proposition~\ref{prop:counterexample} shows neither can be replaced by
``finite Poincar\'e complex'' alone; any proof using only finite complexes and
Poincar\'e duality must fail.
\item \textbf{Odd torsion.} The corresponding statement for odd torsion is
Fowler's theorem \cite{W2}; the present argument adapts the involution case
($p=2$), where the fixed-point and normal-bundle analysis replaces the
Euler-characteristic counting available in nicer settings.
\item \textbf{Role of Novikov.} Only injectivity of the rational assembly map
for $\pi$ is needed, which holds for lattices in semisimple groups via their
proper actions on the nonpositively curved symmetric space.
\item \textbf{Honest accounting.} The delicate points are
\eqref{eq:deg-one}\,--\,\eqref{eq:fixed}: identifying the symmetric signature with
the fundamental class under assembly, and the periodicity/Leray--Hirsch
computation that pushes the fixed-set class up to top dimension. These are the
substance of the surgery-theoretic argument and are stated here at the level of
their structure.
\end{enumerate}

\begin{thebibliography}{9}
\bibitem{Borel} A.~Borel.
  \emph{Seminar on Transformation Groups}. Princeton University Press, 1960.
\bibitem{W1} S.~Weinberger.
  Group actions and higher signatures II.
  \emph{Comm. Pure Appl. Math.}, 1987.
\bibitem{W2} S.~Weinberger.
  \emph{Variations on a Theorem of Borel}.
  Cambridge University Press, 2022.
\end{thebibliography}

\endgroup
% ===== End q7_solution.tex =====

% ===== Begin q8_solution.tex =====
\clearpage
\section*{Problem 8}
\addcontentsline{toc}{section}{Problem 8}
\begingroup
\hypersetup{colorlinks=true,
            linkcolor=blue!50!black,
            citecolor=blue!50!black,
            urlcolor=blue!50!black}
\newcommand{\qeightR}{\mathbb{R}}
\newcommand{\qeightSm}{\mathscr S}
\newcommand{\qeightLdual}{L^{\vee}}
\newcommand{\qeightdd}{\mathrm d}
\newcommand{\qeightLag}{\Lambda}
\newcommand{\qeightHess}{\operatorname{Hess}}
\newtheorem{qeighttheorem}{Theorem}
\newtheorem{qeightproposition}{Proposition}
\newtheorem{qeightlemma}{Lemma}
\newtheorem{qeightcorollary}{Corollary}
\theoremstyle{definition}
\newtheorem{qeightdefinition}{Definition}
\title{Smoothing Polyhedral Lagrangian Surfaces}
\author{}
\date{}
\maketitle


\section*{Statement and Answer}

Equip $\qeightR^4$ with coordinates $(q_1,q_2,p_1,p_2)$ and the standard symplectic
form $\omega=\qeightdd p_1\wedge \qeightdd q_1+\qeightdd p_2\wedge \qeightdd q_2$. A \emph{polyhedral
Lagrangian surface} $K\subset\qeightR^4$ is a finite polyhedral complex, all of whose
faces are Lagrangian, that is a topological submanifold of $\qeightR^4$.

\medskip
\noindent\textbf{Theorem.} \emph{If $K$ is a polyhedral Lagrangian surface in
which exactly $4$ faces meet at every vertex, then $K$ admits a Lagrangian
smoothing: a Hamiltonian isotopy $K_t$ of smooth Lagrangian submanifolds
parametrized by $t\in(0,1]$, extending to a topological isotopy on $[0,1]$ with
$K_0=K$.}

\medskip
\noindent The answer is \textbf{YES}.

\medskip
The proof has a local part (a normal form near each vertex and edge, from which
a smooth structure on $K$ is read off) and a global part (a fibration of
Lagrangian planes transverse to $K$, with respect to which all sufficiently
$C^1$-small ``smoothing functions'' yield genuine graphical Lagrangians). The
$4$-faces hypothesis enters in the local normal form (Lemma~\ref{lem:linear}),
where four Lagrangian quadrants force a standard symplectic basis; it is not a
Maslov/parity condition.

\section*{1. Local model at a vertex}

\begin{qeightlemma}\label{lem:linear}
Let $\qeightR^2\to\qeightR^4$ be an embedding that is linear on each of the four quadrants,
with Lagrangian image $\Sigma$, and assume $\Sigma$ is not contained in a single
plane. Then there is a linear symplectic transformation of $\qeightR^4$ carrying
$\Sigma$ to the product of the union of the positive coordinate axes in
$\qeightR^2_{p_1q_1}$ with that in $\qeightR^2_{p_2q_2}$.
\end{qeightlemma}

\begin{proof}
Let $(v_1,v_2,u_1,u_2)$ be the tangent vectors at the origin to the four edges of
$\Sigma$, ordered so that cyclically adjacent vectors span the faces. The
pairings $\omega(v_i,u_i)$ cannot vanish: if some $\omega(v_i,u_i)=0$, then
$\omega$ would vanish identically on a $3$-dimensional subspace of $\qeightR^4$, which
is impossible for a symplectic form. After swapping $v_i\leftrightarrow u_i$ if
needed we may take both pairings positive, and after rescaling we may take
$\omega(v_1,u_1)=\omega(v_2,u_2)=1$ with the cross pairings vanishing (the faces
are Lagrangian). Then $(v_1,v_2,u_1,u_2)$ is a standard symplectic basis, and the
linear map $\partial_{p_i}\mapsto v_i$, $\partial_{q_i}\mapsto u_i$ is the
required symplectomorphism, carrying $\Sigma$ to the stated product of
positive-axis unions.
\end{proof}

In the plane $\qeightR^2_{pq}$ the symplectic pairing projects the union of the two
positive axes homeomorphically onto the dual of the antidiagonal line $p=q$.
Taking the product and applying Lemma~\ref{lem:linear}:

\begin{qeightcorollary}\label{cor:Ldual}
There is a linear Lagrangian plane $L\subset\qeightR^4$ such that the symplectic
pairing $\qeightR^4\to \qeightLdual$ restricts to a homeomorphism $\Sigma\to \qeightLdual$.
\end{qeightcorollary}

This equips $\Sigma$ with a smooth structure (pulled back from $\qeightLdual$), which
we fix henceforth. Given $L$, call a Lagrangian $\qeightLag\subset\qeightR^4$
\emph{graphical} if the pairing restricts to a diffeomorphism $\qeightLag\cong\qeightLdual$.

\section*{2. Smoothing functions}

Choose a Lagrangian splitting $\qeightR^4\cong L\oplus\qeightLdual$. Each of the four
quadrant-faces of $\Sigma$ is graphical over $\qeightLdual$, hence the graph of the
differential of a quadratic form; collecting these gives quadratic forms
$\{q_{ij}\}_{i,j\in\pm}$ on $\qeightLdual$, agreeing to first order along the edges.
Via $\Sigma\cong\qeightLdual$ (Corollary~\ref{cor:Ldual}) write $q_\Sigma$ for the
resulting $C^1$ function on $\Sigma$ that is $q_{ij}$ on each face.

\begin{qeightdefinition}\label{def:smoothing}
The space $\qeightSm(\Sigma)$ of \emph{smoothing functions} is the set of $C^1$
functions $f:\Sigma\to\qeightR$ such that $f+q_\Sigma$ is $C^\infty$.
\end{qeightdefinition}

$\qeightSm(\Sigma)$ is invariant under adding smooth functions, and depends only on
$L$, not on the splitting (a different splitting changes $q_{ij}$ by a global
quadratic form, i.e.\ a smooth function). To a smoothing function $f$ we
associate the Lagrangian $\qeightLag_{\qeightdd f}\subset\qeightR^4$ that is, piecewise, the graph
of $\qeightdd f$ over each face — equivalently the graph of $\qeightdd(f+q_\Sigma)$ in the
splitting.

\begin{qeightlemma}\label{lem:bijection}
$f\mapsto \qeightLag_{\qeightdd f}$ is a bijection between graphical Lagrangians and
smoothing functions on $\Sigma$ modulo additive constants.
\end{qeightlemma}

\begin{proof}
In the splitting, $\qeightLag_{\qeightdd f}$ is the graph of $\qeightdd(f+q_\Sigma)$ as a function
on $\qeightLdual$; graphical Lagrangians over $\qeightLdual$ are exactly graphs of
differentials of smooth functions, and $f+q_\Sigma$ is smooth precisely when
$f\in\qeightSm(\Sigma)$. The construction does not use the splitting, so the bijection
depends only on $L$.
\end{proof}

\noindent The same statements hold verbatim for an edge: a pair of Lagrangian
half-planes meeting along a line, with the analogue of
Corollary~\ref{cor:Ldual} provided by the next lemma.

\begin{qeightlemma}\label{lem:edge}
If $\Sigma$ is a pair of linear Lagrangian half-planes meeting along a line
$\ell$, the space of Lagrangian planes $L$ for which the pairing
$\Sigma\to\qeightLdual$ is a homeomorphism is contractible.
\end{qeightlemma}

\begin{proof}
Up to affine linear symplectomorphism, $\Sigma$ is the product of the real axis
in one $\qeightR^2$-factor with the union of the positive axes in another. The pairing
$\Sigma\to\qeightLdual$ is a homeomorphism iff $L$ is transverse to both half-planes,
i.e.\ the symplectic reduction of $L$ along $\ell$ is a line in $\ell^\perp/\ell\cong\qeightR^2$
meeting the interior of the positive quadrant — a contractible condition. The
space of Lagrangian lifts of such a line is itself contractible (two lifts differ
by the graph of a map $\ell'\to\ell$), giving the claim.
\end{proof}

\section*{3. Existence of small smoothing functions}

The crux is that $\qeightSm(\Sigma)$ is \emph{not} invariant under rescaling, so the
existence of small smoothing functions is not automatic.

\begin{qeightlemma}\label{lem:small}
$\qeightSm(\Sigma)$ contains functions of arbitrarily small $C^1$-norm.
\end{qeightlemma}

\begin{proof}
Fix a partition of unity $\sum_\sigma \chi_\sigma=1$ on $\Sigma$ indexed by the
strata, with $\chi_\sigma$ supported near $\sigma$ and $\equiv1$ on $\sigma$ away
from its boundary, of bounded $C^k$ norms. Let $\chi_\sigma^\varepsilon$ be the
composition of $\chi_\sigma$ with dilation by $1/\varepsilon$; these are
uniformly bounded with $\|\chi_\sigma^\varepsilon\|_{C^1}=O(1/\varepsilon)$.

Choose a Lagrangian plane $\Lambda_\sigma$ containing each stratum $\sigma$ and
transverse to $L$, with associated smoothing function $f_\sigma$. The first-order
agreement of the faces along shared edges means $f_\sigma$ and $f_{\sigma'}$ agree
to first order along $\sigma\cap\sigma'$. Set $f^\varepsilon:=\sum_\sigma
\chi_\sigma^\varepsilon f_\sigma$. The partition-of-unity structure makes
$f^\varepsilon$ a smoothing function, and although $\|\nabla\chi_\sigma^\varepsilon\|$
grows like $1/\varepsilon$, it is supported where $f_\sigma-f_{\sigma'}=O(\varepsilon^2)$;
by the product rule the $C^1$-norm of $f^\varepsilon$ is therefore $O(\varepsilon)$.
\end{proof}

\section*{4. Global construction: a conormal fibration}

To globalize, we replace the single plane $L$ by a smoothly varying family.

\begin{qeightdefinition}\label{def:fibration}
A \emph{conormal fibration dual to $K$} is a smoothly varying family $L_z$ of
affine Lagrangian planes, $z\in K$, such that: (i) near each vertex the $L_z$ are
translates of a single plane as in Corollary~\ref{cor:Ldual}; (ii) $L_z$ varies
smoothly along edges and along faces; and (iii) in a collar of each edge the
planes in the normal direction are translates of those along the edge, with
\emph{matched} inward normal fields on the two adjacent faces (opposite vectors
under the identification $T_z\sigma\cong\qeightLdual_z\cong T_z\sigma'$).
\end{qeightdefinition}

\begin{qeightlemma}\label{lem:fibexists}
$K$ admits a conormal fibration which near each vertex agrees with the local
choice of Corollary~\ref{cor:Ldual}.
\end{qeightlemma}

\begin{proof}
Lemma~\ref{lem:edge} extends the vertex choices over the edges (the space of
admissible $L_z$ along an edge is contractible). Choosing a normal vector field
to one face along each edge determines matched normals on the other; extending to
the interior of the faces is unobstructed because the space of Lagrangian planes
transverse to a given one is contractible.
\end{proof}

The fibration determines an open neighbourhood $\nu K$ of $K$ in which the fibres
$L_z$ are disjoint, hence a projection $\nu K\to K$. A Lagrangian contained in
$\nu K$ is \emph{graphical} (with respect to $L_z$) if its projection to $K$ is a
diffeomorphism. The local correspondence of Lemma~\ref{lem:bijection} globalizes:

\begin{qeightlemma}\label{lem:graph-global}
Every graphical Lagrangian with respect to $\{L_z\}$ is the graph of a smoothing
function on $K$; and every smoothing function of sufficiently small differential
defines a graphical Lagrangian.
\end{qeightlemma}

\begin{proof}
The correspondence is local on $K$. At a point $z$, a plane $\qeightLdual_z$ transverse
to $L_z$ stays transverse to nearby fibres, so a neighbourhood of $z$ in $\nu K$
is modelled on the conormal bundle of $\qeightLdual_z$ (Weinstein's tubular
neighbourhood theorem), and Lagrangians there are graphs of closed $1$-forms,
i.e.\ differentials of smoothing functions.
\end{proof}

\begin{qeightlemma}\label{lem:small-global}
There exist smoothing functions for $K$ of arbitrarily small $C^1$-norm.
\end{qeightlemma}

\begin{proof}
Take a partition of unity $\sum_\alpha\rho_\alpha=1$ on $K$ indexed by strata,
with $\rho_\alpha$ supported in the open star of $\alpha$. Lemma~\ref{lem:small}
gives smoothing functions $f_\alpha$ of arbitrarily small $C^1$-norm on each open
star; then $\sum_\alpha\rho_\alpha f_\alpha$ is a global smoothing function of
small $C^1$-norm.
\end{proof}

\section*{5. Assembly of the smoothing}

\begin{proof}[Proof of the Theorem]
By Lemma~\ref{lem:fibexists} fix a conormal fibration $\{L_z\}$ with disjoint
fibres on a neighbourhood $\nu K$. By Lemma~\ref{lem:small-global} choose a
sequence of smoothing functions $f_i\to 0$ in $C^1$ with $\qeightdd f_i$ small enough
that, by Lemma~\ref{lem:graph-global}, each $\qeightLag_{\qeightdd f_i}\subset\nu K$ is a
smooth embedded graphical Lagrangian. These $\qeightLag_{\qeightdd f_i}$ are all isotopic to
$K$ by a piecewise-smooth isotopy and converge to $K$ in the Hausdorff topology
as $i\to\infty$ (equivalently as the $C^1$-norm $\to0$).

Because $\qeightLag_{\qeightdd f_i}$ and $\qeightLag_{\qeightdd f_{i+1}}$ are graphs of differentials of
smooth functions over the common base $K$ (via the fibration), the linear
interpolation of those functions produces a smooth Hamiltonian path of graphical
Lagrangians connecting them. Concatenating these paths and reparametrizing yields
a Hamiltonian isotopy $K_t$ ($t\in(0,1]$) of smooth Lagrangians with $K_t\to K$
as $t\to0$, extending to a topological isotopy on $[0,1]$ with $K_0=K$. This is
the desired Lagrangian smoothing.
\end{proof}

\section*{Role of the $4$-faces-per-vertex hypothesis}

The hypothesis is used exactly once, in Lemma~\ref{lem:linear}: four Lagrangian
quadrants meeting at a vertex, with edge tangents $(v_1,v_2,u_1,u_2)$, force
$\omega(v_i,u_i)\ne0$ and hence a standard symplectic basis, which is what makes
the symplectic pairing $\Sigma\to\qeightLdual$ a homeomorphism
(Corollary~\ref{cor:Ldual}) and supplies the smooth structure used throughout.
With a different number of faces the local model changes and this normal form is
unavailable; the mechanism is the linear symplectic geometry of the vertex, not a
Maslov-class parity condition.

\section*{Remarks}

\begin{enumerate}\setlength\itemsep{4pt}
\item \textbf{Why smoothing functions, not Legendrian fillings.} The smoothing is
produced directly as graphs over $K$ via the conormal fibration, so no
holomorphic-curve or contact-fillability input is needed; the only analytic
ingredients are Weinstein's neighbourhood theorem and a partition-of-unity
estimate.
\item \textbf{The rescaling subtlety.} $\qeightSm(\Sigma)$ is closed under adding
smooth functions but not under scalar multiplication, so small smoothing
functions must be built by hand (Lemma~\ref{lem:small}); this is the one place a
genuine estimate is required.
\item \textbf{Matched normals.} Property (iii) of Definition~\ref{def:fibration}
ensures the two smooth structures on $K$ — from Corollary~\ref{cor:Ldual} and from
the edge collars — agree, so the global graph construction is consistent across
edges.
\item \textbf{Generalization.} The argument is local-to-global and uses only
linear symplectic normal forms plus partition-of-unity gluing, so it adapts to
other ambient symplectic manifolds once the vertex normal form is established.
\end{enumerate}

\begin{thebibliography}{9}
\bibitem{Hirsch} M.~W.~Hirsch.
  \emph{Differential Topology}. Graduate Texts in Mathematics, Springer, 1976.
\bibitem{McDuffSalamon} D.~McDuff and D.~Salamon.
  \emph{Introduction to Symplectic Topology}, 3rd ed.,
  Oxford University Press, 2017.
\end{thebibliography}

\endgroup
% ===== End q8_solution.tex =====

% ===== Begin q9_solution.tex =====
\clearpage
\section*{Problem 9}
\addcontentsline{toc}{section}{Problem 9}
\begingroup
\hypersetup{colorlinks=true,
            linkcolor=blue!50!black,
            citecolor=blue!50!black,
            urlcolor=blue!50!black}
\newcommand{\qnineR}{\mathbb{R}}
\newcommand{\qninesgn}{\operatorname{sgn}}
\newcommand{\qninerank}{\operatorname{rank}}
\newcommand{\qninemlrank}{\operatorname{mlrank}}
\newcommand{\qninebm}[1]{\mathbf{#1}}
\newcommand{\qninebQ}{\mathbf{Q}}
\newcommand{\qninebC}{\mathbf{C}}
\newcommand{\qninebA}{\mathbf{A}}
\newcommand{\qninebF}{\mathbf{F}}
\newcommand{\qnineDmat}{D}
\newtheorem{qninetheorem}{Theorem}
\newtheorem{qninelemma}{Lemma}
\newtheorem{qnineproposition}{Proposition}
\newtheorem{qninecorollary}{Corollary}
\title{Algebraic Relations on Determinantal Tensors}
\author{}
\date{}
\maketitle


\section*{Setup and Statement}

For $n\ge 5$, let $A^{(1)},\dots,A^{(n)}\in\qnineR^{3\times4}$ be Zariski-generic.
For each ordered $4$-tuple $(\alpha,\beta,\gamma,\delta)\in[n]^4$ define
$Q^{(\alpha\beta\gamma\delta)}\in\qnineR^{3\times3\times3\times3}$ by
\[
  Q^{(\alpha\beta\gamma\delta)}_{ijk\ell}
   \;=\; \det\!\bigl[A^{(\alpha)}(i,:);\,A^{(\beta)}(j,:);\,A^{(\gamma)}(k,:);\,A^{(\delta)}(\ell,:)\bigr],
   \qquad 1\le i,j,k,\ell\le 3,
\]
the $4\times4$ determinant of the matrix formed by stacking one row from each
of the four matrices.

We seek a polynomial map $\qninebF:\qnineR^{81n^4}\to\qnineR^{N}$ such that:
\begin{enumerate}\setlength\itemsep{2pt}
\item $\qninebF$ does not depend on $A^{(1)},\dots,A^{(n)}$;
\item the degrees of the coordinate functions of $\qninebF$ do not depend on $n$;
\item for $\lambda\in\qnineR^{n\times n\times n\times n}$ with $\lambda_{\alpha\beta\gamma\delta}\ne0$
      precisely on the non-identical tuples,
      $\qninebF\bigl(\lambda_{\alpha\beta\gamma\delta}Q^{(\alpha\beta\gamma\delta)}\bigr)=0$
      \emph{iff} there exist $u,v,w,x\in(\qnineR^{*})^n$ with
      $\lambda_{\alpha\beta\gamma\delta}=u_\alpha v_\beta w_\gamma x_\delta$ for
      all non-identical $(\alpha,\beta,\gamma,\delta)$.
\end{enumerate}

\medskip
\noindent\textbf{Theorem.} \emph{Such a map exists.} Assemble the tensors into
a single $3n\times3n\times3n\times3n$ block tensor $\qninebQ$ whose
$(\alpha,\beta,\gamma,\delta)$ block is $Q^{(\alpha\beta\gamma\delta)}$, and let
$\qninebF$ send the input to the $5\times5$ minors of the four
$3n\times27n^{3}$ mode flattenings of $\qninebQ$. Each coordinate of $\qninebF$ is a
$5\times5$ determinant, hence degree $5$ independent of $n$, and $\qninebF$
satisfies the rank-one characterization.

\medskip
\noindent The answer is \textbf{YES}.

\section*{The key identity: a Tucker decomposition of \texorpdfstring{$\qninebQ$}{Q}}

Index the modes of $\qninebQ$ by $[n]\times[3]$, so that the
$\bigl((\alpha,i),(\beta,j),(\gamma,k),(\delta,\ell)\bigr)$ entry of $\qninebQ$ is
$Q^{(\alpha\beta\gamma\delta)}_{ijk\ell}$.

\begin{qninelemma}\label{lem:tucker}
$\qninebQ$ admits the Tucker decomposition
\begin{equation}\label{eq:tucker}
  \qninebQ \;=\; \qninebC\times_1\qninebA\times_2\qninebA\times_3\qninebA\times_4\qninebA,
\end{equation}
where $\qninebA=[A^{(1)};\dots;A^{(n)}]\in\qnineR^{3n\times4}$ is the vertical
concatenation of the data matrices, and $\qninebC\in\qnineR^{4\times4\times4\times4}$ is
the sign tensor
\[
  C_{abcd}=\begin{cases}\qninesgn(abcd) & a,b,c,d\in[4]\text{ distinct},\\ 0 & \text{otherwise.}\end{cases}
\]
\end{qninelemma}

\begin{proof}
For $(\alpha,i),(\beta,j),(\gamma,k),(\delta,\ell)\in[n]\times[3]$, the
definition of the Tucker product gives
\[
  (\qninebC\times_1\qninebA\times_2\qninebA\times_3\qninebA\times_4\qninebA)_{(\alpha,i),(\beta,j),(\gamma,k),(\delta,\ell)}
  =\sum_{a,b,c,d\in[4]}C_{abcd}\,\qninebA_{(\alpha,i),a}\qninebA_{(\beta,j),b}\qninebA_{(\gamma,k),c}\qninebA_{(\delta,\ell),d}.
\]
Because $\qninebA$ is the vertical concatenation, $\qninebA_{(\alpha,i),a}=A^{(\alpha)}_{ia}$,
and the sum over distinct $a,b,c,d$ weighted by $\qninesgn(abcd)$ is exactly the
Leibniz expansion of the $4\times4$ determinant of the rows
$A^{(\alpha)}(i,:),A^{(\beta)}(j,:),A^{(\gamma)}(k,:),A^{(\delta)}(\ell,:)$.
This equals $Q^{(\alpha\beta\gamma\delta)}_{ijk\ell}=\qninebQ_{(\alpha,i),(\beta,j),(\gamma,k),(\delta,\ell)}$.
\end{proof}

\noindent Equation \eqref{eq:tucker} shows the multilinear rank of $\qninebQ$ is at
most $(4,4,4,4)$: each mode-$m$ flattening factors through the $4$-column matrix
$\qninebA$, so its column space (resp.\ row space) has dimension at most $4$. In
particular \textbf{every $5\times5$ minor of every flattening of $\qninebQ$ vanishes},
and the same holds for any tensor of multilinear rank $\le(4,4,4,4)$. This is
why $\qninebF$, the collection of $5\times5$ minors of the four flattenings, is the
right object: it tests membership in the variety of multilinear rank
$\le(4,4,4,4)$ tensors.

\section*{Definition of \texorpdfstring{$\qninebF$}{F}}

Given input tensors $\{T^{(\alpha\beta\gamma\delta)}\in\qnineR^{3\times3\times3\times3}\}$,
assemble the $3n\times3n\times3n\times3n$ block tensor $\mathbf T$ with blocks
$T^{(\alpha\beta\gamma\delta)}$. Let $\mathbf T_{(m)}\in\qnineR^{3n\times27n^{3}}$ be
its mode-$m$ flattening, $m=1,2,3,4$. Define
\[
  \qninebF(\{T^{(\alpha\beta\gamma\delta)}\})
   \;:=\;\bigl(\text{all }5\times5\text{ minors of }\mathbf T_{(1)},\mathbf T_{(2)},\mathbf T_{(3)},\mathbf T_{(4)}\bigr).
\]
Each coordinate is a $5\times5$ determinant in the entries of the input, hence a
degree-$5$ polynomial in those entries, independent of $n$; and $\qninebF$ visibly
does not reference $A^{(1)},\dots,A^{(n)}$. Properties (1) and (2) hold. We
write $\lambda\odot_b\qninebQ$ for the blockwise scaling whose
$(\alpha,\beta,\gamma,\delta)$ block is $\lambda_{\alpha\beta\gamma\delta}Q^{(\alpha\beta\gamma\delta)}$.
Property (3) is the content of the next two sections.

\section*{The ``if'' direction}

\begin{qnineproposition}\label{prop:if}
If $\lambda_{\alpha\beta\gamma\delta}=u_\alpha v_\beta w_\gamma x_\delta$ on the
non-identical tuples (and $0$ on the diagonal $\alpha=\beta=\gamma=\delta$), then
$\qninebF(\lambda\odot_b\qninebQ)=0$.
\end{qnineproposition}

\begin{proof}
First, the diagonal blocks of $\qninebQ$ vanish: $Q^{(\alpha\alpha\alpha\alpha)}_{ijk\ell}$
is a determinant with rows drawn from the single $3\times4$ matrix $A^{(\alpha)}$,
and any $4\times4$ determinant of vectors lying in a $3$-dimensional row space is
$0$. Hence the value of $\lambda$ on the diagonal is irrelevant and
$\lambda\odot_b\qninebQ=(u\otimes v\otimes w\otimes x)\odot_b\qninebQ$.

Blockwise scaling by a rank-one tensor with nonzero entries is the same as a
Tucker multiplication by invertible diagonal matrices:
\[
  (u\otimes v\otimes w\otimes x)\odot_b\qninebQ
   \;=\;\qninebQ\times_1 D_u\times_2 D_v\times_3 D_w\times_4 D_x,
\]
where $D_u\in\qnineR^{3n\times3n}$ is the diagonal matrix repeating each $u_\alpha$
three times (once per row of the $\alpha$ block), and similarly $D_v,D_w,D_x$.
Multiplying \eqref{eq:tucker} on each mode by an invertible matrix preserves the
multilinear rank, so $\lambda\odot_b\qninebQ$ still has multilinear rank
$\le(4,4,4,4)$, and all its $5\times5$ flattening minors vanish. Thus
$\qninebF(\lambda\odot_b\qninebQ)=0$.
\end{proof}

\section*{The ``only if'' direction}

Assume $\lambda$ has nonzero entries exactly off the diagonal and
$\qninebF(\lambda\odot_b\qninebQ)=0$, i.e.\ $\lambda\odot_b\qninebQ$ has multilinear rank
$\le(4,4,4,4)$. We show $\lambda$ is rank-one off the diagonal.

\subsection*{Normalization}

Replacing $\lambda$ by its entrywise product with $\bar u\otimes\bar v\otimes\bar w\otimes\bar x$,
where e.g.\ $\bar u_1=1$ and $\bar u_\alpha=\lambda_{\alpha111}^{-1}$ for
$\alpha\ge2$ (and $\bar v,\bar w,\bar x$ defined symmetrically on the other
modes), preserves both the multilinear rank of $\lambda\odot_b\qninebQ$ and the
property of being rank-one off-diagonal. Hence we may assume
\begin{equation}\label{eq:normalize}
  \lambda_{\alpha111}=\lambda_{1\beta11}=\lambda_{11\gamma1}=\lambda_{111\delta}=1
  \qquad\text{for all }\alpha,\beta,\gamma,\delta\in\{2,\dots,n\}.
\end{equation}
We will show there is a constant $c\in\qnineR^{*}$ with
\begin{equation}\label{eq:goal}
  \lambda_{\alpha\beta\gamma\delta}=c^{\,m-1},
  \qquad m:=\#\{\text{indices among }\alpha,\beta,\gamma,\delta\text{ that are }\ne 1\},
\end{equation}
for all non-identical tuples. Then $u=v=w=(1,c,\dots,c)$ and
$x=(c^{-1},1,\dots,1)$ give $\lambda_{\alpha\beta\gamma\delta}=u_\alpha v_\beta w_\gamma x_\delta$,
establishing rank one.

\subsection*{Reading off $5\times5$ submatrix determinants}

The multilinear-rank hypothesis says every $5\times5$ minor of every flattening
of $\lambda\odot_b\qninebQ$ vanishes. We extract three families of such minors. Write
$(\lambda\odot_b\qninebQ)_{(1)}\in\qnineR^{3n\times27n^3}$ for the mode-$1$ flattening;
rows are indexed by $(\alpha,i)\in[n]\times[3]$ and columns by
$((\beta,j),(\gamma,k),(\delta,\ell))$.

\medskip
\noindent\textbf{Step 1 (two non-unit indices).}
Choose the $5\times5$ submatrix with rows $(1,1),(1,2),(1,3),(\alpha,1),(\alpha,2)$
and a fixed list of five columns chosen so that, viewed as a polynomial in the
two variables $\lambda_{\alpha\alpha11}$ and $\lambda_{\alpha1\beta1}$, the
determinant has degree at most $2$; it vanishes when $\lambda_{\alpha1\beta1}=0$,
and also when $\lambda_{\alpha1\beta1}=\lambda_{\alpha\alpha11}$ (in the latter
case the submatrix becomes a submatrix of the unscaled $\qninebQ_{(1)}$ after row
operations, hence rank-deficient by Lemma~\ref{lem:tucker}). Therefore this
determinant equals
\[
  s\cdot\lambda_{\alpha1\beta1}\bigl(\lambda_{\alpha1\beta1}-\lambda_{\alpha\alpha11}\bigr),
\]
with $s=s(A^{(1)},A^{(\alpha)},A^{(\beta)})$ a polynomial in the data. A single
generic numerical instance (it suffices to take $\alpha=\beta$, the most
specialized case) shows $s\not\equiv0$, so $s\ne0$ Zariski-generically. The
minor vanishes, forcing $\lambda_{\alpha1\beta1}=\lambda_{\alpha\alpha11}$.
Applying the same argument to all modewise permutations $\pi$ (using
$\pi\cdot\qninebQ=\qninesgn(\pi)\qninebQ$, so the permuted flattening has the same rank bound)
gives $\lambda_{\pi(\alpha1\beta1)}=\lambda_{\pi(\alpha\alpha11)}$. Combining the
specializations $\alpha=\beta$ and $\alpha\ne\beta$ yields a single constant
\[
  c:=\lambda_{\pi(\alpha\beta11)}\quad\text{for all }\alpha,\beta\in\{2,\dots,n\}\text{ and all }\pi,
\]
i.e.\ every $\lambda$-entry with exactly two non-unit indices equals $c$.

\medskip
\noindent\textbf{Step 2 (one non-unit index).}
Take the $5\times5$ submatrix with rows $(1,1),(1,2),(1,3),(\alpha,1),(\alpha,2)$
and columns chosen so the determinant, viewed as a polynomial in $c$ and
$\lambda_{\alpha\beta\gamma1}$, has degree $\le3$, vanishes when $c=0$ (two rows
become dependent), and vanishes when $c^{2}=\lambda_{\alpha\beta\gamma1}$ (the
submatrix reduces to one of $\qninebQ_{(1)}$ after row/column operations). Hence the
determinant is a data-dependent scalar times $c\,(c^{2}-\lambda_{\alpha\beta\gamma1})$.
Generic non-vanishing of the scalar (checked on one instance with
$\alpha=\beta=\gamma$) and the rank hypothesis give
$\lambda_{\alpha\beta\gamma1}=c^{2}$, and by symmetry
$\lambda_{\pi(\alpha\beta\gamma1)}=c^{2}$ for all $\pi$. Thus every $\lambda$-entry
with exactly one non-unit index equals $c^{2}$.

\medskip
\noindent\textbf{Step 3 (no non-unit index).}
Take the $5\times5$ submatrix with rows $(1,1),(1,2),(1,3),(\alpha,1),(\alpha,2)$
and columns producing a determinant equal to a data scalar times
$c\,(c^{3}-\lambda_{\alpha\beta\gamma\delta})$, where now
$\alpha,\beta,\gamma,\delta\in\{2,\dots,n\}$ are not all equal (the construction
parallels Steps 1–2; the most specialized instance $\alpha=\beta=\gamma$ certifies
non-vanishing). The rank hypothesis forces $\lambda_{\alpha\beta\gamma\delta}=c^{3}$,
and by symmetry $\lambda_{\pi(\alpha\beta\gamma\delta)}=c^{3}$.

\medskip
Steps 1–3 establish \eqref{eq:goal}: a $\lambda$-entry with $m$ non-unit indices
equals $c^{\,m-1}$ (using the normalization \eqref{eq:normalize} for $m=1$). As
noted, this is exactly the rank-one factorization
$\lambda_{\alpha\beta\gamma\delta}=u_\alpha v_\beta w_\gamma x_\delta$ with
$u=v=w=(1,c,\dots,c)$ and $x=(c^{-1},1,\dots,1)$. This proves the ``only if''
direction and completes the proof of the Theorem.\qed

\section*{Degree, dimension, and independence}

\textbf{Degree.} Each coordinate of $\qninebF$ is a $5\times5$ determinant in the
entries of the input, so has degree $5$, independent of $n$. \checkmark

\textbf{Independence of the data.} $\qninebF$ is defined purely as minors of the
flattenings of the input block tensor; it never refers to $A^{(1)},\dots,A^{(n)}$. \checkmark

\textbf{Number of coordinates.} $N$ is the number of $5\times5$ minors of four
$3n\times27n^{3}$ matrices, which grows polynomially in $n$; only the
\emph{degree} is required to be bounded, which it is. \checkmark

\section*{Remarks}

\begin{enumerate}\setlength\itemsep{4pt}
\item \textbf{Why $5\times5$ and not $2\times2$.} The obstruction is multilinear
rank $(4,4,4,4)$: a flattening has rank $\le4$, so the vanishing of all
$5\times5$ minors is the exact algebraic condition. Smaller minors do not
characterize the variety, and larger ones are redundant.
\item \textbf{Role of genericity.} Genericity of the $A^{(i)}$ enters only in the
``only if'' direction, to certify that the data-dependent scalars $s$ in
Steps 1–3 are nonzero; each is verified by exhibiting one numerical instance,
after which Zariski-genericity does the rest.
\item \textbf{Diagonal blocks vanish.} $Q^{(\alpha\alpha\alpha\alpha)}=0$ because
four rows of a $3\times4$ matrix are linearly dependent; this is what lets the
off-diagonal rank-one structure of $\lambda$ be detected without the diagonal
interfering.
\item \textbf{Fourth-order analogue.} The construction is the order-$4$ analogue
of the corresponding third-order determinantal-variety result; the block-tensor
flattening and the multilinear-rank bound are the common mechanism.
\end{enumerate}

\begin{thebibliography}{9}
\bibitem{Landsberg} J.~M.~Landsberg.
  \emph{Tensors: Geometry and Applications}. AMS, 2012.
\bibitem{KoldaBader09} T.~G.~Kolda and B.~W.~Bader.
  \emph{Tensor decompositions and applications}.
  SIAM Review \textbf{51} (2009), 455--500.
\bibitem{Sturmfels} B.~Sturmfels.
  \emph{Algorithms in Invariant Theory}, 2nd ed., Springer, 2008.
\end{thebibliography}

\endgroup
% ===== End q9_solution.tex =====

% ===== Begin q10_solution.tex =====
\clearpage
\section*{Problem 10}
\addcontentsline{toc}{section}{Problem 10}
\begingroup
\hypersetup{colorlinks=true,
            linkcolor=blue!50!black,
            citecolor=blue!50!black,
            urlcolor=blue!50!black}
\newcommand{\qtenR}{\mathbb{R}}
\newcommand{\qtenvecop}{\operatorname{vec}}
\newcommand{\qtendiag}{\operatorname{diag}}
\newcommand{\qtenTr}{\operatorname{Tr}}
\newcommand{\qtenKR}{\odot}            % Khatri-Rao
\newcommand{\qtenbigKR}{\bigodot}      % big Khatri-Rao
\newcommand{\qtenhad}{\ast}            % Hadamard / elementwise
\newcommand{\qtenbF}{\bar F}
\newcommand{\qtenbB}{\bar B}
\newcommand{\qtenbW}{\bar W}
\newcommand{\qtenbD}{\bar D}
\newcommand{\qtenbd}{\bar d}
\newcommand{\qtenKhat}{\widehat K}
\newcommand{\qtenZhat}{\widehat Z}
\newcommand{\qtenHhat}{\widehat H}
\newtheorem{qtentheorem}{Theorem}
\newtheorem{qtenlemma}{Lemma}
\newtheorem{qtenproposition}{Proposition}
\title{Preconditioned CG for an RKHS-Constrained CP Subproblem}
\author{}
\date{}
\maketitle


\section*{Setup recap}

We solve the linear system
\begin{equation}\label{eq:sys}
\bigl[\,(Z\otimes K)^{\!\top} S\,S^{\!\top}(Z\otimes K) + \lambda(I_r\otimes K)\bigr]\,\qtenvecop(W)
 \;=\; (I_r\otimes K)\,\qtenvecop(B),
\end{equation}
where $K\in\qtenR^{n\times n}$ is the PSD RKHS kernel matrix on mode $k$,
$Z\in\qtenR^{M\times r}$ is the Khatri--Rao product of the other factor matrices,
$S\in\qtenR^{N\times q}$ is the sampling/selection matrix, $B=TZ\in\qtenR^{n\times r}$
is the MTTKRP, $\lambda>0$ is a Tikhonov regularizer, and the unknown is
$W\in\qtenR^{n\times r}$. We have $n,r,d<q\ll N$ and $n\ll M$, with
$N=\prod_i n_i$ and $M=N/n$. The objective is to avoid any $O(N)$ or $O(M)$
work.

\medskip
\noindent\textbf{Notation.} Let $\Omega\subseteq[N]$, $|\Omega|=q$, index the
observed entries of the vectorized mode-$k$ unfolding; $S^{\!\top}x=x_\Omega$,
$S^{\!\top}S=I_q$, and $SS^{\!\top}=\qtendiag(\mathbf 1_\Omega)$. A vectorized index
$\ell\in[N]$ corresponds to a pair $(i,m)$ with $i\in[n]$ (mode-$k$ index) and
$m\in[M]$ (multi-index of the other modes). Write the observations as triples
$(i_t,m_t,\tau_t)$ for $t=1,\dots,q$.

The system matrix in \eqref{eq:sys},
\[
  H \;:=\; (Z\otimes K)^{\!\top}SS^{\!\top}(Z\otimes K)+\lambda(I_r\otimes K),
\]
is symmetric and PSD (a Gram term plus $\lambda$ times the PSD matrix
$I_r\otimes K$), but it is only positive \emph{semi}definite when $K$ is, and
even when $K\succ0$ it can be poorly conditioned. Our plan is: (i) symmetrize
and add a small shift $\rho$ for definiteness; (ii) diagonalize $K$ to expose a
diagonal preconditioner; (iii) run PCG with matrix-free products built from
three Kronecker-row identities.

\section*{1. Direct symmetric method (baseline)}

Adding a shift $\rho>0$ to ensure positive definiteness, the direct method
forms
\[
  H_\rho \;=\; F^{\!\top}F + \lambda(I_r\otimes K)+\rho I_{rn},
  \qquad F:=S^{\!\top}(Z\otimes K)\in\qtenR^{q\times rn},
\]
and solves $H_\rho\qtenvecop(W)=\qtenvecop(KB)$ by Cholesky. Row $t$ of $F$ is the
Kronecker product
\[
  F(t,:)\;=\;\widehat Z(t,:)\otimes \widehat K(t,:),
  \qquad
  \widehat Z(t,:):=Z(m_t,:),\quad \widehat K(t,:):=K(i_t,:),
\]
so $F^{\!\top}F=\sum_t \bigl(\widehat Z(t,:)\otimes\widehat K(t,:)\bigr)^{\!\top}
\bigl(\widehat Z(t,:)\otimes\widehat K(t,:)\bigr)$. Each rank-one update costs
$O(n^2r^2)$, giving $O(qn^2r^2)$ to form $F^{\!\top}F$, plus $O(n^3r^3)$ for
Cholesky.

\medskip
\noindent\textbf{Total (direct):} $O(qrd+qn^2r^2+n^3r^3)$; storage $O(n^2r^2)$.

\begin{algorithm}[H]
\caption{Direct symmetric solve for the RKHS mode-$k$ subproblem}
\begin{algorithmic}[1]
\Require triples $\{(i_t,m_t,\tau_t)\}_{t=1}^q$, factors $\{A_i\}_{i\ne k}$,
         kernel $K$, regularizer $\lambda$, shift $\rho$.
\Ensure $W\in\qtenR^{n\times r}$.
\State Build sampled Khatri--Rao rows
      $\widehat Z(t,:)\gets\qtenbigKR_{i\ne k}A_i(\mathrm{idx}_i(m_t),:)$ for $t=1,\dots,q$.
\State Build sampled kernel rows $\widehat K(t,:)\gets K(i_t,:)$.
\State Form $F\in\qtenR^{q\times rn}$ row-wise:
      $F(t,:)\gets \widehat Z(t,:)\otimes \widehat K(t,:)$.
\State Assemble
      $H_\rho\gets F^{\!\top}F+\lambda(I_r\otimes K)+\rho I_{rn}$.
\State Accumulate $B(i,:)\gets\sum_{t:i_t=i}\tau_t\,\widehat Z(t,:)$.
\State Form RHS $g\gets\qtenvecop(KB)$.
\State Solve $H_\rho\,\qtenvecop(W)=g$ by Cholesky (or symmetric LDL$^\top$).
\State \Return $W$.
\end{algorithmic}
\end{algorithm}

\section*{2. Transforming the system}

We exploit the eigendecomposition $K=UDU^{\!\top}$ (with $U$ orthogonal, $D$
diagonal) to factor $(I_r\otimes U)$ out of \eqref{eq:sys}. Writing
$\qtenbW:=U^{\!\top}W$ and $\qtenbB:=DU^{\!\top}B$,
\begin{equation}\label{eq:transformed}
  \Bigl[\,\underbrace{(Z\otimes UD)^{\!\top}S}_{\qtenbF^{\!\top}}\,
          \underbrace{S^{\!\top}(Z\otimes UD)}_{\qtenbF}
        + \lambda(I_r\otimes D) + \rho I_{rn}\Bigr]\,\qtenvecop(\qtenbW)
   \;=\; \qtenvecop(\qtenbB).
\end{equation}
We recover $W=U\qtenbW$ after solving. We cannot fold $D$ into $\qtenbW$ because $D$
may be singular (or indefinite after a shift), but $\lambda(I_r\otimes D)$ is
now \emph{diagonal}, which is the crucial gain: the regularization and the
shift together contribute a known diagonal, and (\S4) the diagonal of
$\qtenbF^{\!\top}\qtenbF$ can be computed exactly and cheaply. Define
\[
  \qtenbF := S^{\!\top}(Z\otimes UD)\in\qtenR^{q\times rn},
  \qquad H_\rho^{\,\prime}:=\qtenbF^{\!\top}\qtenbF+\lambda(I_r\otimes D)+\rho I_{rn}.
\]
Row $t$ of $\qtenbF$ is again a Kronecker product of two rows,
\begin{equation}\label{eq:Fbar-rows}
  \qtenbF(t,:)=\widehat Z(t,:)\otimes \widehat H(t,:),
  \qquad \widehat H(t,:):=H(i_t,:),\quad H:=UD\in\qtenR^{n\times n},
\end{equation}
so all three lemmas of \S3 apply with $A=\widehat Z$ and $B=\widehat H$.

\section*{3. Three Kronecker-row identities}

The matrix-free products and the preconditioner rest on the following
identities. Let $A\in\qtenR^{q\times r}$ and $B\in\qtenR^{q\times n}$, and define
$C\in\qtenR^{q\times rn}$ row-wise by
\begin{equation}\label{eq:Crows}
  C(t,:) \;=\; A(t,:)\otimes B(t,:),\qquad t=1,\dots,q.
\end{equation}
Under the column-major vec convention, an index $\ell\in[rn]$ corresponds to
$(i,j)$ with $i\in[n]$, $j\in[r]$ via $\ell=i+(j-1)n$; then
$C_{t\ell}=B_{ti}A_{tj}$ and $(\qtenvecop X)_\ell=X_{ij}$.

\begin{qtenlemma}[forward product]\label{lem:Cx}
For $X\in\qtenR^{n\times r}$ and $x=\qtenvecop(X)$,
\[
  Cx \;=\; \bigl(A\qtenhad (BX)\bigr)\,\mathbf 1_r,
\]
where $\qtenhad$ is the elementwise product and $\mathbf 1_r$ the all-ones vector.
\end{qtenlemma}
\begin{proof}
$(Cx)_t=\sum_{\ell}C_{t\ell}x_\ell=\sum_{j=1}^r\sum_{i=1}^n B_{ti}X_{ij}A_{tj}
 =\sum_{j=1}^r (BX)_{tj}A_{tj}$, which is row $t$ of $(A\qtenhad BX)\mathbf 1_r$.
\end{proof}

\begin{qtenlemma}[adjoint product]\label{lem:CTv}
For $v\in\qtenR^q$, $\;C^{\!\top}v=\qtenvecop\!\bigl(B^{\!\top}\qtendiag(v)A\bigr).$
\end{qtenlemma}
\begin{proof}
With $\ell=i+(j-1)n$,
$(C^{\!\top}v)_\ell=\sum_t C_{t\ell}v_t=\sum_t B_{ti}A_{tj}v_t
 =(B^{\!\top}\qtendiag(v)A)_{ij}=(\qtenvecop(B^{\!\top}\qtendiag(v)A))_\ell.$
\end{proof}

\begin{qtenlemma}[diagonal of the Gram]\label{lem:diag}
$\;\qtendiag(C^{\!\top}C)=\qtenvecop\!\bigl((B\qtenhad B)^{\!\top}(A\qtenhad A)\bigr).$
\end{qtenlemma}
\begin{proof}
$(C^{\!\top}C)_{\ell\ell}=\sum_t C_{t\ell}^2=\sum_t B_{ti}^2A_{tj}^2
 =\bigl((B\qtenhad B)^{\!\top}(A\qtenhad A)\bigr)_{ij}.$
\end{proof}

\noindent Lemma~\ref{lem:Cx} computes $Cx$ in $O(q(r+n))$; Lemma~\ref{lem:CTv}
computes $C^{\!\top}v$ in $O(q(r+n))$; Lemma~\ref{lem:diag} computes the full
diagonal of $C^{\!\top}C$ in $O(q(r^2+n^2))$ — all without forming $C$, which
would need $O(qrn)$ storage.

\section*{4. PCG: matrix-free products and preconditioner}

We never form $\qtenbF$, $Z$, $S$, or $H_\rho^{\,\prime}$. PCG needs the
right-hand side, the product $w\mapsto H_\rho^{\,\prime}w$, and a
preconditioner.

\subsection*{4.1 The matrix--vector product}

Apply the row factorization \eqref{eq:Fbar-rows} with $A=\widehat Z$,
$B=\widehat H$ where $\widehat H(t,:)=H(i_t,:)$, $H=UD$. For
$x=\qtenvecop(X)$, $X\in\qtenR^{n\times r}$:
\begin{itemize}\setlength\itemsep{2pt}
\item By Lemma~\ref{lem:Cx}, $\qtenbF x=(\widehat Z\qtenhad(\widehat H X))\mathbf 1_r\in\qtenR^q$.
      Forming $\widehat H X$ uses only the $q$ needed rows of $HX$; in practice
      compute $HX=U(DX)\in\qtenR^{n\times r}$ once at $O(n^2r)$ and gather rows
      $i_t$.
\item By Lemma~\ref{lem:CTv}, $\qtenbF^{\!\top}(\qtenbF x)
      =\qtenvecop\!\bigl(\widehat H^{\!\top}\qtendiag(\qtenbF x)\,\widehat Z\bigr)$,
      an $O(q(r+n))$ scatter.
\end{itemize}
Hence
\begin{equation}\label{eq:matvec}
  H_\rho^{\,\prime}x
   = \qtenvecop\!\Bigl(\widehat H^{\!\top}\,\qtendiag\!\bigl((\widehat Z\qtenhad \widehat H X)\mathbf 1_r\bigr)\,\widehat Z\Bigr)
     + \lambda\,DX + \rho X,
\end{equation}
computed in $O(n^2r+q(r+n))$ per iteration (the $\lambda(I_r\otimes D)$ and
$\rho I_{rn}$ terms are diagonal, hence $O(nr)$).

\subsection*{4.2 The diagonal preconditioner}

Because the transformed regularizer is diagonal, the natural Jacobi
preconditioner is \emph{exact and cheap}. Define
\[
  \qtenbD := \qtendiag\!\bigl(\qtendiag(\qtenbF^{\!\top}\qtenbF)\bigr)+\lambda(I_r\otimes D)+\rho I_{rn},
\]
whose diagonal $\qtenbd=\qtendiag(\qtenbD)$ is, by Lemma~\ref{lem:diag},
\begin{equation}\label{eq:precond}
  \qtenbd \;=\; \qtenvecop\!\bigl((\widehat H\qtenhad \widehat H)^{\!\top}(\widehat Z\qtenhad \widehat Z)\bigr)
        \;+\;\lambda\,(\mathbf 1_r\otimes\qtendiag(D))\;+\;\rho\,\mathbf 1_{rn}.
\end{equation}
This costs $O(q(r^2+n^2))$ once, and each application
$\qtenbD^{-1}z=z/\qtenbd$ costs $O(nr)$. Unlike a block-Tikhonov choice, this is the
true diagonal of the operator, so it captures the per-coordinate scale of the
sampling term exactly rather than approximately. The conditioning of
$\qtenbD^{-1}H_\rho^{\,\prime}$ is governed by how far $\qtenbF^{\!\top}\qtenbF$ is from
diagonal; in the well-sampled, well-conditioned-$Z$ regime it is close, and
PCG converges in few iterations. We make no concentration claim beyond the
standard CG bound $\ell=O(\sqrt{\kappa(\qtenbD^{-1}H_\rho^{\,\prime})}\,\log(1/\varepsilon))$.

\subsection*{4.3 Algorithm}

\begin{algorithm}[H]
\caption{PCG for the transformed RKHS mode-$k$ subproblem}
\begin{algorithmic}[1]
\Require triples $\{(i_t,m_t,\tau_t)\}_{t=1}^q$, factors $\{A_i\}_{i\ne k}$,
         kernel $K$, regularizer $\lambda$, shift $\rho$, tolerance $\varepsilon$.
\Ensure $W\in\qtenR^{n\times r}$.
\State \textbf{Setup:}
\State \quad Eigendecompose $K=UDU^{\!\top}$;\; set $H\gets UD$. \Comment{$O(n^3)$, once}
\State \quad Sampled Khatri--Rao rows $\widehat Z(t,:)\gets\qtenbigKR_{i\ne k}A_i(\mathrm{idx}_i(m_t),:)$. \Comment{$O(qrd)$}
\State \quad Gather $\widehat H(t,:)\gets H(i_t,:)$ for $t=1,\dots,q$. \Comment{$O(qn)$}
\State \quad $B(i,:)\gets\sum_{t:i_t=i}\tau_t\,\widehat Z(t,:)$;\; RHS $g\gets\qtenvecop(DU^{\!\top}B)$. \Comment{$O(qr+n^2r)$}
\State \quad Precompute $\qtenbd$ via \eqref{eq:precond}. \Comment{$O(q(r^2+n^2))$}
\State \textbf{PCG:} $w_0\gets0$, $r_0\gets g$, $z_0\gets r_0/\qtenbd$, $p_0\gets z_0$.
\For{$\ell=0,1,2,\dots$ until $\|r_\ell\|\le\varepsilon\|g\|$}
  \State $a_\ell\gets H_\rho^{\,\prime}p_\ell$ via \eqref{eq:matvec}. \Comment{$O(n^2r+q(r+n))$}
  \State $\alpha_\ell\gets\langle r_\ell,z_\ell\rangle/\langle p_\ell,a_\ell\rangle$.
  \State $w_{\ell+1}\gets w_\ell+\alpha_\ell p_\ell$;\quad $r_{\ell+1}\gets r_\ell-\alpha_\ell a_\ell$.
  \State $z_{\ell+1}\gets r_{\ell+1}/\qtenbd$;\quad
         $\beta_\ell\gets\langle r_{\ell+1},z_{\ell+1}\rangle/\langle r_\ell,z_\ell\rangle$.
  \State $p_{\ell+1}\gets z_{\ell+1}+\beta_\ell p_\ell$.
\EndFor
\State \Return $W\gets U\cdot\mathrm{unvec}(w_\ell)$. \Comment{$O(n^2r)$}
\end{algorithmic}
\end{algorithm}

\section*{5. Complexity analysis}

\textbf{Setup (PCG):}
$O\bigl(n^3 + qrd + q(r^2+n^2) + n^2r\bigr)$. The eigendecomposition of $K$ is
$O(n^3)$, performed once before the outer alternating iterations; nothing of
size $M$ or $N$ is ever formed.

\textbf{Per iteration:} $O\bigl(n^2r + q(r+n)\bigr)$. The $HX=U(DX)$ products
($O(n^2r)$), the Lemma~\ref{lem:Cx}/\ref{lem:CTv} gather--scatter
($O(q(r+n))$), and the diagonal solve ($O(nr)$).

\textbf{Iterations:}
$\ell=O\bigl(\sqrt{\kappa(\qtenbD^{-1}H_\rho^{\,\prime})}\,\log(1/\varepsilon)\bigr)$.

\textbf{Total (PCG):}
\[
  O\bigl(n^3 + qrd + q(r^2+n^2) + \ell\cdot(n^2r+q(r+n))\bigr).
\]

\textbf{Comparison.} Direct costs $O(qn^2r^2+n^3r^3)$. PCG removes the
Kronecker $r^2$ blow-up in the iteration cost (down to $O(n^2r)$) and avoids
the $O(n^3r^3)$ Cholesky, at the price of $\ell$ matrix-free products. Memory
drops from $O(n^2r^2)$ (storing $H$) to $O(n^2+nr+q(r+n))$ (kernel factors,
iterate, sampled rows).

\section*{6. Why no object of size \texorpdfstring{$M$}{M} or \texorpdfstring{$N$}{N} appears}

\begin{enumerate}\setlength\itemsep{4pt}
\item \emph{Kronecker--vec.} $(Z\otimes K)\qtenvecop(W)=\qtenvecop(KWZ^{\!\top})$, but
$KWZ^{\!\top}\in\qtenR^{n\times M}$ is never formed; only the $q$ entries selected
by $\Omega$ are needed, and Lemma~\ref{lem:Cx} produces exactly those from the
$q$ sampled rows.
\item \emph{Sampling avoids $M$.} The forward map yields a length-$q$ vector of
inner products; the adjoint (Lemma~\ref{lem:CTv}) scatters back to an
$n\times r$ matrix. Both are $O(q(r+n))$.
\item \emph{$Z$ is never formed.} Only the $q$ rows $\widehat Z(t,:)$ that are
used appear, each an $r$-fold Hadamard product of factor rows at $O(rd)$ — the
sampled Khatri--Rao trick.
\item \emph{Diagonalizing $K$.} The transform $K=UDU^{\!\top}$ turns the
regularizer into a diagonal, enabling the exact diagonal preconditioner
\eqref{eq:precond} via Lemma~\ref{lem:diag}; the only $O(n^3)$ cost is the
one-time eigendecomposition.
\item \emph{$K$ is small.} With $n\ll M$, the $O(n^2r)$ multiplications and
$O(n^3)$ factorization are affordable.
\end{enumerate}

\section*{7. Definiteness}

$H$ is symmetric PSD: $F^{\!\top}F\succeq0$ and $\lambda(I_r\otimes K)\succeq0$.
When $K\succ0$, $H\succ0$. When $K$ is only PSD, the shift $\rho>0$ (equivalently
$K\to K+\eta I$) restores strict positive definiteness so that CG is
well-defined; in the transformed coordinates this appears as the
$\rho I_{rn}$ term and a possibly-shifted diagonal $D$. The transform is
orthogonal, so it preserves the spectrum of the regularized operator and the
recovered $W=U\qtenbW$ solves the original system.

\section*{8. Summary}

\begin{center}\small
\begin{tabular}{@{}lll@{}}\toprule
 & Direct & PCG (transformed) \\\midrule
Setup & $O(qrd+n^2r)$ & $O(n^3+qrd+q(r^2+n^2))$ \\
Per iter / solve & $O(qn^2r^2+n^3r^3)$ & $O(n^2r+q(r+n))$ \\
Iterations & 1 (Cholesky) & $\ell=O(\sqrt{\kappa}\,\log\tfrac1\varepsilon)$ \\
Preconditioner & --- & exact diagonal (Lemma~\ref{lem:diag}) \\
Memory & $O(n^2r^2)$ & $O(n^2+nr+q(r+n))$ \\
Forms $H$? & yes & no \\
Touches $M$ or $N$? & no & no \\
\bottomrule\end{tabular}\end{center}

\medskip
\noindent\textbf{Conclusion.} Diagonalizing the kernel converts the
regularizer into a diagonal, after which the three Kronecker-row identities
give matrix-free products and an \emph{exact} diagonal (Jacobi) preconditioner.
The result solves the RKHS mode-$k$ subproblem in
$O(\ell\,(n^2r+q(r+n)))$ work per iteration plus $O(n^3)$ one-time setup,
without forming the $nr\times nr$ system matrix or any object of size
$\Omega(M)$ or $\Omega(N)$.

\begin{thebibliography}{9}
\bibitem{KoldaBader09} T.~G.~Kolda and B.~W.~Bader.
  \emph{Tensor decompositions and applications}.
  SIAM Review \textbf{51} (2009), 455--500.
\bibitem{Saad03} Y.~Saad.
  \emph{Iterative Methods for Sparse Linear Systems}, 2nd ed., SIAM, 2003.
\end{thebibliography}

\endgroup
% ===== End q10_solution.tex =====

\end{document}
